\documentclass[11pt]{article}
% Shared preamble for the paper and its supplement (pdflatex).
\PassOptionsToPackage{hyphens}{url}
\usepackage[T1]{fontenc}
\usepackage[utf8]{inputenc}
\usepackage{lmodern}
\usepackage[letterpaper,margin=1in]{geometry}
\usepackage{setspace}
\setstretch{1.12}
\usepackage{amsmath,amssymb,amsthm,mathtools}
\usepackage{booktabs,array,tabularx,longtable}
\usepackage{enumitem}
\usepackage{xcolor}
\usepackage{microtype}
\usepackage[authoryear,round]{natbib}
\usepackage{xr-hyper}
\usepackage[hidelinks]{hyperref}
\urlstyle{same}
\setlength{\emergencystretch}{3em}
\setlist{itemsep=2pt,topsep=4pt}
\newcolumntype{L}[1]{>{\raggedright\arraybackslash\hspace{0pt}}p{#1}}
\newcolumntype{Y}{>{\raggedright\arraybackslash}X}

% ---------------- notation ----------------
\newcommand{\Bset}{\mathfrak B}
\newcommand{\Kspace}{\mathcal K}
\newcommand{\ThI}{\Theta_I}
\newcommand{\ThX}{\Theta_X}
\newcommand{\ThO}{\Theta_O}
\newcommand{\E}{\mathbb E}
\newcommand{\R}{\mathbb R}
\newcommand{\conv}{\operatorname{conv}}
\newcommand{\supp}{\operatorname{supp}}
\newcommand{\vertx}{\operatorname{vert}}
\newcommand{\affdim}{\operatorname{affdim}}
\newcommand{\Var}{\operatorname{Var}}
\newcommand{\Cov}{\operatorname{Cov}}
\newcommand{\sym}{\operatorname{sym}}
\newcommand{\skw}{\operatorname{skew}}
\newcommand{\diag}{\operatorname{diag}}
\newcommand{\hull}{\operatorname{hull}}
\newcommand{\sgn}{\operatorname{sign}}
\newcommand{\one}{\mathbf 1}
\newcommand{\Bt}{\tilde B}
\newcommand{\DD}{\mathrm{DD}}
\newcommand{\CI}{\mathrm{CI}}
\newcommand{\ER}{\mathrm{ER}}
\newcommand{\NM}{\mathrm{NM}}
\newcommand{\pkg}[1]{\textsf{#1}}
\newcommand{\code}[1]{\texttt{#1}}

% ---------------- result environments ----------------
\theoremstyle{plain}
\newtheorem{theorem}{Theorem}
\newtheorem{proposition}{Proposition}
\newtheorem{lemma}{Lemma}
\newtheorem{corollary}{Corollary}
\newtheorem{conjecture}{Conjecture}
\theoremstyle{definition}
\newtheorem{assumption}{Assumption}
\newtheorem{definition}{Definition}
\newtheorem{procedure}{Procedure}
\theoremstyle{remark}
\newtheorem{remark}{Remark}
\newtheorem{example}{Example}

% plain-language reading and sketch of the argument
\newenvironment{reading}{\par\smallskip\noindent\textit{Reading.}\ }{\par\smallskip}
\newenvironment{sketch}{\par\smallskip\noindent\textit{Sketch of the argument.}\ }{\par\smallskip}
% status tag printed after a result's title
\newcommand{\status}[1]{\textup{[\textsf{#1}]}}

\externaldocument[S-]{eiPartialID_supplement}

\title{Partial Identification for R$\times$C Ecological Inference:\\
Identified Sets, Inference and Software}
\author{[Author]}
\date{November 2026}

\begin{document}
\maketitle

\begin{abstract}
\noindent Ecological inference asks how groups behaved when only group shares
and outcome shares are observed for each unit. Point estimates require strong
assumptions that the data cannot check. This paper develops the alternative:
state an assumption, derive the set of values of the target that are consistent
with the data and that assumption, and report that set with honest sampling
uncertainty. We extend the $2\times2$ bounds of Jiang, King, Schmaltz and Tanner
under linear contextual effects to general R$\times$C tables. The data identify
the symmetric part of the contextual-effect matrices and leave a skew part of
dimension $(C-1)R(R-1)/2$ unidentified. We characterise the sharp identified
set, show that the computable linear-programming bound is sharp for the
regression function, and give closed forms for two rows. For units nested in
districts we derive a ladder of nested assumptions, from Duncan--Davis bounds
through district-specific and shared contextual effects to Goodman constancy.
We add sign and contextual-reinforcement beliefs from monotone ecological
inference as linear restrictions, and give exact sharpness and redundancy
results for the pooled model. Because each endpoint of the bound is the minimum
of many estimated candidate limits, the plug-in bound is biased inward; we give
confidence intervals with proved asymptotic coverage, a finite-sample interval
valid in every regime, tests of the model, and a sensitivity analysis. The
default interval of the accompanying R package, a CLR-type procedure with
simulated candidates, is stated as an algorithm whose validity is conjectured.
Simulations and a benchmark with known truth from New Zealand elections show
where each interval is informative.
\end{abstract}

\noindent\textbf{Keywords:} ecological inference; partial identification;
intersection bounds; linear programming; contextual effects; monotonicity.

\tableofcontents
\clearpage

% Section 1
\section{Introduction}\label{sec:intro}

\subsection{The problem}

Ecological inference (EI) asks how groups behaved when only aggregates are
observed. A unit is a polling place, a precinct, a county or a riding. For
each unit we see the shares of its members in $R$ groups (how its electorate
voted in an earlier election, or its ethnic composition) and the shares of its
members in $C$ outcome categories (how it voted in a later election, or whether
it turned out). We want the latent $R\times C$ table of rates: the share of the
members of group $r$ who chose outcome $c$, unit by unit and, above all, for
the electorate as a whole. Transition rates between elections, turnout by
group, and racially polarised voting are the standard examples.

The data link groups and outcomes only through an accounting identity: in each
unit, the outcome shares are the group shares weighted by the unknown rates.
A unit has $RC$ unknown rates, but its margins supply only $C-1$ free
equations. Without further structure the rates are not point-identified, and
neither is their electorate-wide average.

The classical answers resolve this in one of two ways. The first assumes enough
to obtain a single number. Goodman's ecological regression assumes that the
rates are the same in every unit~\citep{goodman1953}. King's method assumes a
truncated bivariate normal distribution for the unit rates~\citep{king1997}, and
hierarchical R$\times$C models make related distributional
assumptions~\citep{rjkt2001,greinerquinn2009}. These assumptions do the
identifying work, and the data cannot check them. When they fail, for example
because group behaviour depends on the composition of the unit, the answer can
be confidently wrong. The second answer assumes nothing beyond the accounting
identity and reports the widest possible range, the Duncan--Davis ``method of
bounds''~\citep{duncan1953}. It is always valid, but often too wide to be
useful.

\subsection{The aim of the bounding approach}

Partial identification takes the middle road~\citep{manski2003,tamer2010,molinari2020}.
It states an assumption, derives the set of values of the target that are
consistent with the data and that assumption (the \emph{identified set}), and
reports that set. When the assumption is weak the set is wide, but it is
credible. When the assumption is strong the set is narrow, but it rests on more.
The researcher can see exactly what each assumption buys.

\citet{jkst2019} (JKST below) carried out this programme for $2\times2$ tables.
Their assumption is \emph{linear contextual effects} (LCE): each group's
expected rate may depend on the composition of the unit, but only linearly. LCE
is much weaker than Goodman's constancy, which is the special case with no
contextual effect. JKST show that under LCE the aggregate regression of the
outcome share on the group share is quadratic, that it identifies all but one
number, and that requiring the implied rates to be probabilities bounds that
number and hence the target.

This paper extends the programme in three directions and sets out what can be
learned under each assumption, with honest uncertainty:
\begin{enumerate}
\item \textbf{General R$\times$C tables.} The unidentified part grows from one
number to a $(C-1)R(R-1)/2$-dimensional set, and the bound becomes a linear
program (LP).
\item \textbf{A ladder of assumptions.} Units are often nested in districts
(precincts in ridings, counties in states). District baselines protect against
aggregation bias across districts, and shared contextual effects carry the
information of the most informative district to all of them. Sign beliefs from
monotone ecological inference~\citep{egghm2025} add further, cheap restrictions.
Each assumption is a rung, and rungs are reported side by side.
\item \textbf{Inference.} Every bound is computed from an estimated regression.
Each endpoint is the tightest of many estimated candidate limits, so the
plug-in bound is biased inward, sometimes badly. Valid confidence intervals must
treat each endpoint as an intersection bound~\citep{clr2013}.
\end{enumerate}

\subsection{Contributions}

The results are of four kinds, and we keep them apart throughout. A result
labelled \emph{proved} has a complete proof in the supplementary appendix. A
\emph{cited} result applies a published theorem whose conditions we verify. A
\emph{conjecture} is stated precisely, supported by checks or simulation, and
not proved. An \emph{algorithm} is a procedure whose validity is not proved;
simulation evidence does not promote it to a theorem.

\paragraph{Identification in R$\times$C.} Under LCE the regression of the outcome
shares on the group shares identifies the symmetric part $Q$ of the
contextual-effect matrices and nothing about the skew part $K$
(Lemma~\ref{lem:decomp}). We characterise the sharp identified set $\ThI$ for
$K$ through conditional support-function inequalities (Theorem~\ref{thm:sharp}),
define two larger sets $\ThX$ and $\ThO$, prove the nesting
$\ThI\subseteq\ThX\subseteq\ThO$ (Proposition~\ref{prop:nesting}), and show that
the outer set $\ThO$ is exactly what the regression function can deliver
(Proposition~\ref{prop:sharpreg}). $\ThO$ is a polytope, and bounds on any linear
target are an LP (Proposition~\ref{prop:lp}). For two groups the LP has a closed
form (Proposition~\ref{prop:tworows}). In $2\times2$ the construction reproduces
JKST's published bounds exactly (Proposition~\ref{prop:jkst}).

\paragraph{Districts and the rung ladder.} Grouping units adds nothing without a
model that links groups (Proposition~\ref{prop:D0}). With such a model, the
sharpness theorem carries over to any mean model that is affine in its unknowns
(Theorem~\ref{thm:affine}), and the rungs form a ladder of nested identified sets
(Theorem~\ref{thm:ladder}). We give a closed form for shared contextual effects
with two groups, which shows how one district with lopsided units and an extreme
baseline tightens the bound everywhere (Proposition~\ref{prop:D3}), and the exact
condition under which the pooled model is valid (Proposition~\ref{prop:D4}).

\paragraph{Signs, reinforcement and monotonicity.} Beliefs about the sign of the
between-group and within-group (contextual) associations, and the belief that the
two reinforce each other, are accounting identities away from data-computable
quantities (Lemma~\ref{lem:M1}). On realised tables they give LP bounds that are
sharp for continuous unit tables (Proposition~\ref{prop:N1}). Under the pooled
linear rung each becomes one linear row in $K$, the restricted bound is sharp for
the regression (Theorem~\ref{thm:R3}), and a face test decides exactly when a
belief tightens a target (Proposition~\ref{prop:R2}). These sign rows are
implemented in the accompanying package as an \emph{experimental} module: the
results are proved for the pooled model only, and sharpness and inference under
district rungs are open. A monotone nonparametric rung gives a valid outer LP
(Proposition~\ref{prop:N2}); its sharp version is a conjecture
(Conjecture~\ref{conj:N2}).

\paragraph{Inference.} Each endpoint is a minimum of finitely many linear
functions of the regression coefficients (Lemma~\ref{lem:minlin}). The plug-in
endpoint has a non-normal limit with negative mean when candidates tie
(Theorem~\ref{thm:limit}, Proposition~\ref{prop:triangular}). We prove asymptotic
coverage for a conservative interval and for projection and fixed-superset
intervals (Theorem~\ref{thm:ci}), give a variance-adaptive allowance for the
realised rate (Proposition~\ref{prop:bd}), tests of the model
(Proposition~\ref{prop:tests}), a finite-sample interval valid in every
asymptotic regime (Theorem~\ref{thm:finite}), and inference with added sign rows
and either/or restrictions (Theorem~\ref{thm:R10}). The CLR-type interval with
simulated candidates, the package default, is narrower in grouped designs, but
its validity is Conjecture~\ref{conj:clr}.

\paragraph{Software and evidence.} The R package \pkg{eiPartialID} implements the
rungs, the intervals and the tests, and prints for each fit which result covers
its intervals. A Monte Carlo study with 1000 replications per design, a benchmark
with known truth from New Zealand elections, and simulations of the sign module
show where each interval is informative and what happens when an assumption is
false.

\subsection{Related work}

The bounds of \citet{duncan1953} are the model-free baseline. \citet{crossmanski2002}
study ecological inference with an instrument that makes some rates constant
across units and obtain a point; LCE is weaker, and the bounds here lie between
their data-alone region and their point. The random-set formulation of sharp
identified sets follows \citet{bmm2011}. The inference builds on intersection
bounds~\citep{clr2013}, directionally differentiable functionals
\citep{shapiro1991,fangsantos2019}, and the distinction between covering the
parameter and covering the identified set~\citep{imbensmanski2004,stoye2009}.
The sign and reinforcement beliefs come from \citet{egghm2025} (EGGHM below);
their identities hold for any bounded outcome and are used here as accounting
identities. The view of the sources of uncertainty follows
\citet{aaiw2023} and \citet{aiz2014}.

\subsection{Plan}

Section~\ref{sec:setting} sets out the data, the targets, the model and the
assumptions, each with a plain-language reading. Section~\ref{sec:rxc} gives
the identification results for R$\times$C tables under pooled LCE.
Section~\ref{sec:districts} adds districts and the rung ladder.
Section~\ref{sec:signs} adds sign, reinforcement and monotonicity beliefs.
Section~\ref{sec:inference} turns the identified sets into confidence intervals
and tests. Section~\ref{sec:package} describes the R package,
Section~\ref{sec:evidence} the numerical evidence and a worked example, and
Section~\ref{sec:open} the open questions. The supplementary appendix contains
the proofs in one notation.

In the main text each formal statement is followed by a plain-language reading
and, for the key results, a sketch of the argument. A sketch is a guide to the
proof, not a substitute for it.

% Section 2
\section{Setting, targets and assumptions}\label{sec:setting}

\subsection{Data and the accounting identity}

There are units $i=1,\dots,n$. Unit $i$ has size $N_i$ (its number of
members), row shares $X_i\in\Delta^{R-1}$ and column shares
$T_i\in\Delta^{C-1}$, where $\Delta^{k-1}=\{x\in\R^k:x\ge0,\ \one^\top x=1\}$
is the probability simplex and $\one$ is the vector of ones. Rows are groups
(or choices in an earlier election), indexed by $r,s=1,\dots,R$; columns are
outcomes, indexed by $c=1,\dots,C$. When units are nested in districts,
$D_i\in\{1,\dots,D^\star\}$ is the district of unit $i$, and $V_i=(X_i,N_i,D_i)$
collects the design variables (without districts, $V_i=(X_i,N_i)$).

The latent rate $\beta_{irc}$ is the share of the row-$r$ members of unit $i$ who
fall in column $c$. Each row of the matrix $\beta_i=[\beta_{irc}]$ lies in the
simplex, so $\beta_i\in\Bset:=(\Delta^{C-1})^R$. The accounting identity is
\begin{equation}\label{eq:acc}
T_{ic}=\sum_{r=1}^R X_{ir}\beta_{irc},\qquad\text{that is,}\qquad T_i=\beta_i^\top X_i .
\end{equation}
For given margins $(x,t)$, the \emph{unit Duncan--Davis polytope}
\[
S(x,t)=\{b\in\Bset:\ b^\top x=t\}
\]
is the set of rate matrices consistent with one unit's margins, and
$h_S(u;x,t)=\max_{b\in S(x,t)}\langle u,b\rangle$ is its support function, where
$\langle u,b\rangle=\sum_{rc}u_{rc}b_{rc}$. With cell masses $n_{rc}=x_rb_{rc}$
the polytope is a transportation polytope, and its single-cell extremes are the
Fréchet--Hoeffding bounds $\max\{0,x_r+t_c-1\}\le n_{rc}\le\min\{x_r,t_c\}$.

\subsection{Three targets}\label{sec:targets}

Three related targets appear. Each result says which one it concerns. Write
$w_{ir}=N_iX_{ir}/\sum_jN_jX_{jr}$ for the share of row-$r$ members who live in
unit $i$, and $m(v)=\E[\beta\mid V=v]$ for the conditional mean of the rates.
\begin{itemize}
\item The \emph{actual} rate $B_{rc}=\sum_iw_{ir}\beta_{irc}$: what the row-$r$
members of this electorate did. It is JKST's target.
\item The \emph{model-implied} rate $\Bt_{rc}=\sum_iw_{ir}m_{rc}(V_i)$: what the
systematic pattern implies they did, at the observed design. It is the
conditional mean of $B$ given the design.
\item The \emph{superpopulation} rate $B^*_{rc}=\E[NX_r\beta_{rc}]/\E[NX_r]$,
used in the population identification results.
\end{itemize}
The actual and model-implied rates differ by a residual term,
$B_{rc}=\Bt_{rc}+E_{rc}$ with $E_{rc}=\sum_iw_{ir}e_{irc}$ and
$e_i=\beta_i-m(V_i)$. A linear target is $\varphi=\sum_{rc}a_{rc}B_{rc}$ for
fixed weights $a_{rc}$; loyalty, defection, demobilisation and the gap between
two groups' rates are examples. All inference in this paper is conditional on
the design and covers $B$ or $\Bt$; inference for $B^*$ would need the
sampling variation of the design and is not pursued.

\begin{reading}
$B$ is what these voters actually did. $\Bt$ is what the systematic pattern
implies they did; it ignores the local luck of each unit. A confidence interval
for $B$ must allow for that luck, so it is wider.
\end{reading}

\subsection{Linear contextual effects in homogeneous form}

Because $\one^\top x=1$ on the simplex, an affine function of $x$ can be written
as a homogeneous linear one: $a+b^\top x=(a\one+b)^\top x$. The pooled LCE model,
called M0, says that for each column $c$ there is a matrix $H_c\in\R^{R\times R}$
with
\begin{equation}\label{eq:lce}
\E[\beta_{irc}\mid X_i,N_i]=e_r^\top H_cX_i\qquad\text{on }\{X_i\in A\},
\end{equation}
where $e_r$ is the $r$-th unit vector and $A\subseteq\Delta^{R-1}$ is the compact
convex region on which the model is assumed. The default region is the
population hull $A=\conv\supp P_X$. Write $H=(H_1,\dots,H_C)$ and
$M_H(x)=[e_r^\top H_cx]_{rc}$ for the implied mean rate matrix at composition
$x$. Split each matrix into its symmetric and skew parts,
\[
H_c=Q_c+K_c,\qquad Q_c=\tfrac12(H_c+H_c^\top),\qquad K_c=\tfrac12(H_c-H_c^\top).
\]
Taking conditional expectations of \eqref{eq:acc}, the regression function is
\[
\tau_c(x):=\E[T_c\mid X=x]=x^\top H_cx=x^\top Q_cx ,
\]
because $x^\top K_cx=0$ for a skew matrix. The regression therefore reveals $Q$
and never sees $K$. The unknown $K=(K_1,\dots,K_C)$ lives in
\[
\Kspace=\Big\{(K_1,\dots,K_C):\ K_c^\top=-K_c,\ \textstyle\sum_cK_c=0\Big\},
\qquad\dim\Kspace=(C-1)\,R(R-1)/2 ,
\]
which is one number in $2\times2$, two in $2\times3$ and six in $3\times3$. We
write $k$ for the free coordinates $k_{c,rs}$ ($c<C$, $r<s$) of $K$, and, for
$R=2$, $\kappa_c=K_{c,12}$.

\paragraph{What $H$, $Q$ and $K$ mean.} $H_{c,rs}$ is the rate at which row-$r$
members choose column $c$ in a hypothetical unit made up entirely of row-$s$
members. $Q_{c,rs}$ is the average of ``$r$-members in pure-$s$ units'' and
``$s$-members in pure-$r$ units''; the data identify it. $K_{c,rs}$ is half the
gap between the two; the data never identify it. This is the R$\times$C version
of JKST's question: is turnout low in a mixed area because group A turns out
less, or because group B turns out less where A is numerous?

\paragraph{Moments.} The superpopulation weights are $\mu_r=\E[NX_r]>0$ and
$\mu_{rs}=\E[NX_rX_s]$. In finite-design statements the same symbols denote the
sample analogues with weights $p_i=N_i/\sum_jN_j$ (Section~\ref{sec:signs}).

\subsection{The model and three identified sets}

Let $P$ be the law of the observables $W=(N,X,T)$ (with $D$ when there are
districts), and $\mathbb Q$ a law of $(W,\beta)$. The model $\mathcal M(H;A)$ is the
set of laws $\mathbb Q$ such that
\begin{itemize}
\item[(S1)] $\beta\in\Bset$ and $T=\beta^\top X$ almost surely (the accounting identity);
\item[(S2)] $\E[\beta\mid V]=M_H(X)$ almost surely on $\{X\in A\}$ (linear contextual
effects, with no dependence on size);
\item[(S3)] $M_H(x)\in\Bset$ for every $x\in A$ (the linear model gives valid
probabilities on all of $A$).
\end{itemize}
Let $Q(P)$ be the unique symmetric matrices with $\E_P[T_c\mid X=x]=x^\top Q_cx$
(Lemma~\ref{lem:decomp} gives existence and uniqueness under the model). Three
sets of values of $K$ are used.
\begin{itemize}
\item The \emph{sharp} identified set
$\ThI(P;A)=\{K\in\Kspace:\ \exists\,\mathbb Q\in\mathcal M(Q(P)+K;A)$ whose
$W$-marginal is $P\}$. It uses the whole law of the data and conditions on
$(X,N)$.
\item $\ThX(P;A)$: the same with (S2) replaced by the weaker
$\E[\beta\mid X]=M_H(X)$. This is the set of JKST's Proposition~1.
\item The \emph{outer} set
$\ThO(P;A)=\{K\in\Kspace:\ (Q_c(P)+K_c)x\ge0$ elementwise for all $c$ and all
$x\in A\}$. It uses only the regression function and the requirement that the
implied rates be probabilities.
\end{itemize}
Proposition~\ref{prop:nesting} shows $\ThI\subseteq\ThX\subseteq\ThO$. Three
hulls also appear: the population hull $\conv\supp P_X$ (the default $A$ in
identification), the hull of the empirical law of the design (exact for
finite-design identification), and the sample hull used as an estimate of $A$.
The sample hull lies inside the population hull, so estimating it can only widen
the outer bound (Proposition~\ref{prop:consistency}(c)). Two conventions hold
throughout. First, $\ThO$ is called sharp only \emph{for the regression function}
(Proposition~\ref{prop:sharpreg}). Second, every ``sharp'' statement is relative
to continuous (real-valued) latent rates, not integer cell counts.

\subsection{Assumptions and what each implies}\label{sec:assumptions}

Each assumption is stated and then read in plain terms: what it rules in or
out, when it is credible, and what goes wrong if it fails.
Table~\ref{tab:assumption-use} shows which rungs use which assumption.

\begin{assumption}[A1: closed population]\label{ass:A1}
$T_i=\beta_i^\top X_i$ exactly, with the same $N_i$ people in both margins.
\end{assumption}
\begin{reading}
Both margins describe the same people: the same electorate in both elections,
or the same eligible population for group and outcome. A1 rules out turnover,
moving, deaths, new registrants and boundary changes. It is credible when both
margins come from one register and the categories include abstention and,
where relevant, ``entered'' or ``left''. If it fails, the Duncan--Davis sets and
the regression identity both shift, and every bound can exclude the truth.
Bounded turnover changes both the unit polytopes and the regression identity,
so it cannot be handled by widening the final interval; it is an open
extension. Every rung, including Duncan--Davis, needs A1.
\end{reading}

\begin{assumption}[A2: linear contextual effects on $A$]\label{ass:A2}
Each group's mean rate is linear in the unit's composition on $A$:
equation~\eqref{eq:lce}.
\end{assumption}
\begin{reading}
A group's behaviour may depend on its surroundings, but only along a straight
line: for example, loyalty rises steadily as the in-party share of the unit
rises. A2 rules out thresholds and tipping (a machine that mobilises only above
50\%) and U-shapes (competitive units mobilising both sides). It is the
\emph{only} identifying restriction of the linear rungs. If it fails, the
quadratic regression is misspecified and the bound can miss the truth by an
amount that does not shrink with more data. Some violations are invisible:
nonlinear departures that cancel in the aggregate leave the data unchanged while
moving the target (Proposition~\ref{prop:undetectable}). A cubic term or an empty
bound can reveal some failures; goodness of fit can never validate A2.
\end{reading}

\begin{assumption}[A3: mean independence from size]\label{ass:A3}
$\E[\beta\mid X,N]=\E[\beta\mid X]$, with $D$ added to both conditioning sets
under the district rungs.
\end{assumption}
\begin{reading}
Large and small units with the same mix behave alike. This matters because the
target weights units by size. If large urban units differ in loyalty from small
rural ones with the same mix, the size-weighted target is biased by
$\Cov(NX_r,\beta_{rc}-m_{rc}(X))/\E[NX_r]$. District baselines help when size
differs by district (Proposition~\ref{prop:D4}). A3 is partly testable through
interactions of $X$ with $\log N$ in the regression.
\end{reading}

\begin{assumption}[A4: support and extrapolation]\label{ass:A4}
$\supp(X)\subseteq A$, $A$ is compact, and A2 holds on all of $A$.
\end{assumption}
\begin{reading}
The straight lines are trusted over the whole region $A$, whose corners (the
most extreme units) set the constraints. A larger $A$ gives tighter bounds but
assumes more: taking the whole simplex assumes that the lines extend to pure
units that may not exist (Proposition~\ref{prop:simplex}). Outlying units at the
corners of the hull can make constraints bind spuriously; trimming the hull is
the check. Probability validity on $A$, condition (S3), is automatic for the
population hull (Lemma~\ref{lem:outer}(c)).
\end{reading}

\begin{assumption}[A5: rank]\label{ass:A5}
For every symmetric $R\times R$ matrix $\Psi$, $x^\top\Psi x=0$ for all
$x\in\supp_AP_X$ implies $\Psi=0$. Here $\supp_AP_X$ is the support of the
restriction of $P_X$ to $A$: the smallest closed set $F$ with $P(X\in A\setminus F)=0$.
\end{assumption}
\begin{reading}
The compositions vary enough to fit the quadratic regression. For two groups it
means at least three distinct compositions. For three groups it needs six
compositions that do not lie on a common conic; any five do, so five are never
enough (Remark~\ref{rem:A5}). A5 fails, or nearly fails, when shares barely vary.
Then $Q$ is weakly identified and the bound is wide and unstable. The condition
number of the design matrix is the check. The restricted support matters only
when $A$ carries less than full mass: for $X$ uniform on $[0,1]$ and
$A=\{1/2\}$, $\supp P_X\cap A=\{1/2\}$ but $\supp_AP_X=\emptyset$.
\end{reading}

\begin{assumption}[A6: source of randomness and independence]\label{ass:A6}
Given the design $\{(N_i,X_i,D_i)\}$ and the district effects, the unit rates
$\beta_i$ are random with mean $m(V_i)$, and the residuals $e_i=\beta_i-m(V_i)$ are
independent across units.
\end{assumption}
\begin{reading}
Every unit is observed, so the uncertainty does not come from sampling units. It
is model-based: given its characteristics, each unit's behaviour has a random
part. A shock common to a district is a fixed feature of this electorate: part
of the mean under the district rungs, a misspecification under the pooled one.
Independence is a substantive assumption. Neighbourhood spillovers within a
district violate it, and then the variances and residual allowances are too
small. Observing every unit does not remove dependence. The pooling test of
Section~\ref{sec:inference} tests mean shifts, not dependence.
\end{reading}

\begin{assumption}[A7: regularity for the normal limit]\label{ass:A7}
A unique active constraint at the optimum of the LP, $\mu_{rs}>0$, and $\ThO$ not
a singleton.
\end{assumption}
\begin{reading}
This is JKST's tie-breaking condition. With it the endpoint of the bound is a
smooth function of the regression coefficients and the delta method applies.
In R$\times$C, and especially with many districts, near-ties are the norm, so the
inference of Section~\ref{sec:inference} does not need it. If it fails, the endpoint
is only directionally differentiable, the delta method loses its justification
and the standard bootstrap is inconsistent (Theorem~\ref{thm:limit}). This is a
failure of the inference, not a bias in identification.
\end{reading}

\begin{assumption}[A8: shared contextual effects across districts]\label{ass:A8}
The contextual effects are common to all districts; only the intercepts vary.
\end{assumption}
\begin{reading}
The way composition shifts behaviour is the same in every district, while each
district may be more or less loyal across the board (parallel lines). A8
produces the intersection gain of the shared-effects rung. It fails if, say, a
strong incumbent changes the contextual gradient in one district, or if floor
and ceiling effects make an additive shift impossible near 0 or 1, which is most
likely in the lopsided districts that tighten the bound most. The bound can then
exclude the truth and still look narrow. The symmetric part of the restriction
is testable (Lemma~\ref{lem:D1}(b)); the shared $K$ is not.
\end{reading}

\begin{assumption}[A9: within-district richness]\label{ass:A9}
Under district-specific effects, A5 within each district. Under shared effects,
each district's support contains $R$ affinely independent compositions and the
stacked design has full column rank (Lemma~\ref{lem:D1}(d)).
\end{assumption}
\begin{reading}
Districts with few, similar units identify almost nothing. Under shared effects
the bar is lower: a district's baseline needs only $R$ affinely independent
compositions. Weak identification per district, combined with a maximum or
minimum over districts, biases plug-in bounds inward.
\end{reading}

\begin{assumption}[A10: asymptotic regime]\label{ass:A10}
One of: (R1) $D^\star$ fixed and the number of units per district $n_d\to\infty$;
(R2) $D^\star\to\infty$ with $n_d$ bounded; (R3) both grow.
\end{assumption}
\begin{reading}
Which limit theory applies depends on the design. It is a modelling choice, not
a fact. With ten polling places in each of ninety ridings (R2), the district
intercepts are incidental parameters and the plug-in intersection is badly too
narrow. Table~\ref{tab:regimes} in Section~\ref{sec:inference} says which result
holds in which regime.
\end{reading}

\paragraph{Assumptions not made.} No distributional form for $\beta$ (unlike
King's truncated normal), no independence between $\beta$ and $X$ (contextual
effects are the point), no stationarity over time, and no point-identifying
nonlinearity.

\begin{table}[t]
\centering\small
\begin{tabularx}{\linewidth}{@{}L{0.33\linewidth}Y@{}}
\toprule
Rung & Assumptions used \\
\midrule
Duncan--Davis (DD) & A1 only. \\
Sign rows on realised tables (N1) & A1; continuous unit tables; no regression, no sampling inference. \\
Pooled LCE (M0) & A1--A5; inference adds A6, and A7 only for the normal limit. \\
District rungs (M2, M1($\delta$), M1, L4, L5, L6) & A1--A4 within districts, A9 in place of A5, A8 for M1 and the rungs below it; inference adds A6 and the regime A10. \\
Monotone means (N2) & A1, A3, and the finite-support convention of Section~\ref{sec:N2}. \\
Sign rows under a linear rung & Those of the rung they are added to. \\
\bottomrule
\end{tabularx}
\caption{Which rungs use which assumption. The rungs are defined in Sections~\ref{sec:districts} and~\ref{sec:signs}.}\label{tab:assumption-use}
\end{table}

% Section 3
\section{Identification in \texorpdfstring{R$\times$C}{RxC} under linear contextual effects}\label{sec:rxc}

This section concerns the pooled model M0 and the superpopulation target $B^*$.
Throughout, assume $\E[N]<\infty$ and $\mu_r>0$ for every $r$.

\subsection{The unit polytope and what the data identify}

\begin{lemma}[The unit polytope]\label{lem:polytope}\status{proved}
For all $x\in\Delta^{R-1}$ and $t\in\Delta^{C-1}$: (a) $S(x,t)\neq\emptyset$;
(b) $S(x,t)$ is compact and convex; (c) $(x,t)\mapsto S(x,t)$ has a closed graph,
is upper hemicontinuous and Borel measurable; (d) $h_S(u;x,t)$ is sublinear and
$\sqrt R$-Lipschitz in $u$, concave in $t$, and Borel in $(u,x,t)$; (e)
$S(x,t)=\{b:\langle u,b\rangle\le h_S(u;x,t)\ \forall u\}$.
\end{lemma}
\begin{reading}
Each unit's set of possible rate tables is never empty (``every group behaves
like the unit as a whole'', $b_{rc}=t_c$, always fits) and it is a well-behaved
convex polytope. These are the regularity facts every later proof uses.
\end{reading}

\begin{lemma}[Identification decomposition]\label{lem:decomp}\status{proved}
Let $\mathbb Q\in\mathcal M(H;A)$ have $W$-marginal $P$, and assume A5. Then
(a) $\E[T_c\mid X=x]=x^\top Q_cx$ on $A$, and $Q_c$ is unique;
(b) $\sum_cH_c=\one\one^\top$, so $\sum_cQ_c=\one\one^\top$ and $\sum_cK_c=0$;
(c) every $K'\in\Kspace$ gives the same regression function, and
$\dim\Kspace=(C-1)R(R-1)/2$.
\end{lemma}
\begin{reading}
The data reveal the symmetric part $Q$ of the contextual-effect matrices and
nothing about the skew part $K$. For each pair of groups and each outcome, the
data identify the sum ``effect of the $s$-share on group $r$ plus effect of the
$r$-share on group $s$'', but not how it splits.
\end{reading}
\begin{sketch}
Conditional expectations of the accounting identity give $\E[T_c\mid X]=X^\top
H_cX$, and a skew matrix vanishes in a quadratic form, so only $Q_c$ appears. A5
makes the quadratic form determine $Q_c$. Row sums of $\beta$ equal one; under A5
this forces $\sum_cH_c=\one\one^\top$.
\end{sketch}

\begin{remark}[A5 in practice]\label{rem:A5}
For $R=2$, A5 holds iff $\supp_AP_X$ contains at least three distinct points. In
general A5 says that the design of the quadratic regression has full rank. For
$R=3$ that needs six support points not on a common conic, and any five points
lie on a conic. Fitting each column by ordinary least squares (OLS) on the
$R(R+1)/2$ quadratic monomials gives $\sum_c\hat Q_c=\one\one^\top$ exactly,
because $\sum_cT_{ic}=1=X_i^\top\one\one^\top X_i$; no constrained regression is
needed.
\end{remark}

\begin{lemma}[The estimand]\label{lem:estimand}\status{proved}
Under (S1)--(S2) with $P(X\in A)=1$,
\[
B^*_{rc}=\frac1{\mu_r}\Big[\sum_sQ_{c,rs}\mu_{rs}+\sum_{s\ne r}K_{c,rs}\mu_{rs}\Big],
\]
which is affine in $K$, with coefficient $\mu_{rs}/\mu_r\ge0$ on the entry $K_{c,rs}$.
\end{lemma}
\begin{reading}
The target is a known number plus a weighted sum of the unknown entries of $K$.
The weight $\mu_{rs}/\mu_r$ measures how much group $r$ lives alongside group
$s$: the more mixed the units, the more the unknown split matters. In pure
units it does not matter at all.
\end{reading}

\subsection{The sharp identified set}

Let $P_{T\mid V}$ be a regular conditional distribution of $T$ given $V$, and
\[
\psi(u;v):=\int h_S(u;x,t)\,P_{T\mid V=v}(dt),
\]
the conditional expectation of the support function at $V=v=(x,n)$.

\begin{theorem}[Characterisation of the sharp set]\label{thm:sharp}\status{proved}
Assume A1, A5 and $H=Q(P)+K$ with $K\in\Kspace$. Then $K\in\ThI(P;A)$ if and
only if (S3) holds and there is one $P_V$-null set $\mathcal N$ such that for
every $v=(x,n)\notin\mathcal N$ with $x\in A$,
\begin{equation}\label{eq:T1}
\langle u,M_H(x)\rangle\le\psi(u;v)\qquad\text{for all }u\in\R^{R\times C}.
\end{equation}
$\ThI(P;A)$ is convex and closed. It is compact when $\affdim(A)=R-1$, which A5
implies.
\end{theorem}
\begin{reading}
A value of $K$ is consistent with the data exactly when, at every composition and
size, the implied mean rate table could be the average of rate tables each of
which fits some unit's margins. Formally, $M_H(x)$ lies in the conditional
selection expectation of the random polytope $S(X,T)$ given $V=v$. This is the
strongest possible use of the data under the model. It needs the whole
conditional distribution of outcomes given composition and size, which in real
data has to be smoothed, so it is not what the package computes. Choosing $u$
and $-u$ along the equality constraints shows that \eqref{eq:T1} also requires
$\E[T\mid V]=\E[T\mid X]$, an observable implication of A3.
\end{reading}
\begin{sketch}
Necessity: each unit's $\beta$ lies in its polytope, so its conditional mean
satisfies the support-function inequality direction by direction. One null set
for all directions comes from a countable dense set of directions and the
Lipschitz continuity of $h_S$ in $u$. Sufficiency: the set of measurable
selections of $S(X,T)$ is weakly compact in $L^1$ (Dunford--Pettis and Mazur),
and so is its image under $\E[\cdot\mid V]$. If the target mean were outside that
image, the Hahn--Banach theorem would separate it by a bounded $V$-measurable
direction, and \eqref{eq:T1} evaluated at that random direction (a measurable
argmax attains the support function) would give a contradiction. No auxiliary
randomisation is needed, because $S(x,t)$ is convex. The random-set
formulation follows \citet{bmm2011}.
\end{sketch}

\begin{corollary}[Sharp bounds for a linear target]\label{cor:sharp}\status{proved}
If $P(X\in A)=1$ and $\ThI\neq\emptyset$, the sharp identified set for
$\varphi=\sum_{rc}a_{rc}B^*_{rc}$ is the interval
$[\min_{K\in\ThI}\varphi(K),\ \max_{K\in\ThI}\varphi(K)]$, with $\varphi(K)$ given
by Lemma~\ref{lem:estimand}.
\end{corollary}
\begin{reading}
The unknown enters the target linearly and the consistent set is convex and
compact, so the answer is always an interval.
\end{reading}

\paragraph{Continuous versus integer rates.} Sharpness is relative to real-valued
latent rates. With integer cell counts the exact feasible set is a subset of
$S(X,T)$. Every vertex of $S(x,t)$ with integer margins is an integer table,
because transportation polytopes are totally unimodular, so the Duncan--Davis
bounds are attained by integer tables. Whether randomising over integer vertices
leaves the sharp set unchanged has only a sketch of an argument, and nothing in
the paper uses it.

\subsection{The outer set: reduction, nesting and sharpness for the regression}

\begin{lemma}[Vertex reduction, boundedness, validity]\label{lem:outer}\status{proved}
(a) If $A$ is a polytope and $\ell$ is linear, $\ell\ge0$ on $A$ iff $\ell(v)\ge0$
at every vertex $v$ of $A$; so the constraints of $\ThO$ may be imposed at
$\vertx(A)$ only. (b) If $\affdim(A)=R-1$, $\ThO(P;A)$ is bounded, hence a
polytope; otherwise it is a closed convex polyhedron, possibly unbounded.
(c) Under (S1)--(S2), if $A\subseteq\conv(\supp_AP_X)$, then (S3) holds
automatically; in particular for $A=\conv\supp P_X$.
\end{lemma}
\begin{reading}
Only the most extreme units matter for the constraints. If the region $A$ is the
range of the data, probability validity of the straight lines is automatic. If
$A$ extrapolates beyond the data, or reaches into a part of the hull where the
model is imposed on no support point, validity there is an extra assumption.
For example, with $A$ a single interior point that carries no mass, (S2)
restricts nothing and $M_H$ need not be valid there; and a single-point $A$ can
make $\ThO$ unbounded. A5 rules out the second case, because it forces $A$ to be
full-dimensional.
\end{reading}

\begin{proposition}[Nesting]\label{prop:nesting}\status{proved}
$\ThI(P;A)\subseteq\ThX(P;A)\subseteq\ThO(P;A)$.
\end{proposition}
\begin{reading}
Using more of the data can only narrow the set. The outer set uses only the
regression function, so it is the widest of the three. The second inclusion is
JKST's ``Jensen'' step; in R$\times$C it needs no Jensen inequality, because
conditional expectations of valid rate tables are valid rate tables.
\end{reading}

\begin{proposition}[$\ThO$ is sharp for the regression function]\label{prop:sharpreg}\status{proved}
Let $\mathcal P(P)$ be the set of laws $P'$ of $(N,X,T')$ with the same law of
$(N,X)$ as $P$ and $\E_{P'}[T'\mid V]=\E_P[T\mid X]$. If $P(X\in A)=1$, then
\[
\ThO(P;A)=\bigcup_{P'\in\mathcal P(P)}\ThI(P';A).
\]
\end{proposition}
\begin{reading}
The outer bound is exactly what can be learned from the regression alone. It is
the sharp set for any data set in which units with the same composition have
identical outcomes. Calling it ``sharp'' without that qualifier would be wrong.
\end{reading}
\begin{sketch}
Given $K\in\ThO$, set every unit's rates equal to the implied mean, $\beta=M_H(X)$.
These are valid rates because $K\in\ThO$, they reproduce the regression, and $K$
is in the sharp set of the resulting law. The reverse inclusion follows from
Proposition~\ref{prop:nesting}, since $\ThO$ depends on $P$ only through $Q(P)$.
\end{sketch}

\begin{proposition}[When the outer set is sharp]\label{prop:IequalsO}\status{proved}
Assume $P(X\in A)=1$. (a) If $T=\E[T\mid X]$ almost surely, then $\ThI=\ThO$.
(b) More generally, suppose $\E[T\mid V]=\E[T\mid X]$ and, for $P_V$-almost every
$v=(x,n)$ and every $u$, $t\mapsto h_S(u;x,t)$ is affine on the convex hull of the
conditional support of $T$ given $V=v$. Then $\ThI=\ThO$. The condition in (b) is
sufficient, not necessary.
\end{proposition}
\begin{reading}
The regression-only bound loses nothing unless units with similar compositions
are so different that some sit at a Duncan--Davis corner and others do not. In
$2\times2$, the condition says that among units with the same composition the
outcome share never lies on both sides of $x$ or of $1-x$, where the unit bounds
have their kinks. In 400 random valid $2\times2$ designs the inclusion
$\ThX\subseteq\ThO$ was strict in 46.
\end{reading}

\subsection{Closed forms and the linear program}

\begin{proposition}[Equivalence with JKST]\label{prop:jkst}\status{proved}
Let $R=C=2$, $X$ the row-1 share, $A=\{(t,1-t):t\in[l,u]\}$ with $0<l<u<1$, and
write $\E[T_1\mid X]=w_0+c_1X+d_1X^2$ and $K=K_{1,12}$. Then (a)
$Q_{1,22}=w_0$, $Q_{1,12}=w_0+c_1/2$, $Q_{1,11}=w_0+c_1+d_1$ and $K=c_1/2-w_1$,
where $w_1$ is JKST's unidentified slope; (b) $\ThO$, written in terms of $w_1$, is
JKST's Proposition~2; (c) with $A=\Delta^1$ it is JKST's Proposition~3; (d) if
$P(l\le X\le u)=1$, $B^*(K)$ is the population version of JKST's equation~(16).
\end{proposition}
\begin{reading}
The R$\times$C machinery contains the published $2\times2$ method exactly, and
the package's original bounds are a special case. Numerically, the homogeneous
form reproduces the original implementation to $10^{-16}$. The resulting bounds
are over the outer set, as are JKST's: sharp for the regression function, not for
the full law.
\end{reading}

\begin{proposition}[The whole simplex]\label{prop:simplex}\status{proved}
With $A=\Delta^{R-1}$, $\ThO=\{K\in\Kspace:\ |K_{c,rs}|\le Q_{c,rs}\ \forall c,\ r<s\}$,
which also requires $Q_{c,rr}\ge0$ for all $c,r$ (a testable condition not involving $K$).
It is non-empty iff every $Q_{c,rs}\ge0$; then each $K_{c,rs}$ ranges over
$[-m_{c,rs},m_{c,rs}]$ with $m_{c,rs}=\min\{Q_{c,rs},1-Q_{c,rs}\}\le\tfrac12$, and
\[
B^*_{rc}\in\frac1{\mu_r}\Big[\sum_sQ_{c,rs}\mu_{rs}\mp\sum_{s\ne r}m_{c,rs}\mu_{rs}\Big].
\]
A negative $Q_{c,rs}$ refutes the model with $A=\Delta^{R-1}$.
\end{proposition}
\begin{reading}
If the straight lines are trusted all the way to pure units, two hypothetical
probabilities with a known average $Q$ can split it by at most $\pm\min(Q,1-Q)$.
This is narrower than the hull-based bound but assumes extrapolation.
\end{reading}

\begin{lemma}[Box intersected with a hyperplane]\label{lem:box}\status{proved}
Let $\ell\le h$ in $\R^C$ and $\mathcal P=\{\kappa:\ell\le\kappa\le h,\ \sum_c\kappa_c=0\}$.
(a) $\mathcal P\neq\emptyset$ iff $\sum_c\ell_c\le0\le\sum_ch_c$.
(b) The projection of $\mathcal P$ on coordinate $c$ is
$\big[\max\{\ell_c,-\sum_{c'\ne c}h_{c'}\},\ \min\{h_c,-\sum_{c'\ne c}\ell_{c'}\}\big]$.
(c) $\max_{\mathcal P}g^\top\kappa$ is attained by the greedy rule: start from
$\kappa=\ell$ and raise coordinates in decreasing order of $g_c$, each up to
$h_c$, until $\sum_c\kappa_c=0$.
\end{lemma}

\begin{proposition}[Two rows, any number of columns]\label{prop:tworows}\status{proved}
Let $R=2$ and $\vertx(A)=\{(t_k,1-t_k)\}_{k=1,2}$ with $0<t_1<t_2<1$. Define
\[
\lambda_c(t)=-Q_{c,12}-Q_{c,11}\tfrac{t}{1-t},\qquad
\rho_c(t)=Q_{c,12}+Q_{c,22}\tfrac{1-t}{t},\qquad
\ell_c=\max_k\lambda_c(t_k),\qquad h_c=\min_k\rho_c(t_k).
\]
Then $\ThO=\{\kappa:\ell\le\kappa\le h,\ \sum_c\kappa_c=0\}$, each $\kappa_c$ is
projected by Lemma~\ref{lem:box}(b), and if $P(X\in A)=1$ the bounds on
$B^*_{1c}=[Q_{c,11}\mu_{11}+(Q_{c,12}+\kappa_c)\mu_{12}]/\mu_1$ (increasing in
$\kappa_c$), on $B^*_{2c}=[(Q_{c,12}-\kappa_c)\mu_{12}+Q_{c,22}\mu_{22}]/\mu_2$
(decreasing in $\kappa_c$), and on any linear target (by the greedy rule) are in
closed form.
\end{proposition}
\begin{reading}
For two groups and any number of outcomes, each outcome's unknown is squeezed
between the tightest lower and upper limits implied by the two extreme units (the
least and the most group-1-heavy), and the unknowns must sum to zero across
outcomes. No LP solver is needed. Group 1's rate rises and group 2's rate falls
as the unknown rises, because the unknown is the share of the observed pattern
attributed to group 1.
\end{reading}

\begin{corollary}[Joint versus collapsed analysis]\label{cor:joint}\status{proved}
Let $R=2$. The collapsed analysis of column $c$ treats the table as two columns,
$c$ and ``not $c$''. (a) The joint interval for $\kappa_c$ lies inside the
collapsed one. (b) They are equal whenever, for each row $r$, the numbers
$Q_{c',rr}$ ($c'\neq c$) share one sign; in particular when every extrapolated
pure-unit rate $H_{c',rr}=Q_{c',rr}$ is non-negative. (c) Strict tightening
therefore requires some negative extrapolated pure-unit rate, so the linear model
is not valid on the whole simplex.
\end{corollary}
\begin{reading}
With two groups, fitting all outcome columns jointly helps a single column only if
the straight lines would give an impossible rate at some pure unit. With three or
more groups the rows are coupled through the shared entries of $K$, and joint
fitting can be strictly tighter even when every extrapolation is valid (strict
in 20 of the $3\times3$ cases checked). A characterisation for $R\ge3$ is open.
Joint fitting is needed anyway for targets that combine columns.
\end{reading}

\begin{proposition}[The general linear program]\label{prop:lp}\status{proved}
Assume $P(X\in A)=1$ and that $A$ is a polytope. In the free coordinates $k$,
with $K_c=\sum_{r<s}k_{c,rs}(e_re_s^\top-e_se_r^\top)$ for $c<C$ and
$K_C=-\sum_{c<C}K_c$,
\[
\ThO=\{k:\ Gk\le g\},
\]
with one row for each (vertex $v$, row $r$, column $c$): $-e_r^\top K_c(k)v\le(Q_cv)_r$.
A linear target is $\varphi(k)=\varphi_0+\omega^\top k$ with
\[
\varphi_0=\sum_{rc}\frac{a_{rc}}{\mu_r}\sum_sQ_{c,rs}\mu_{rs},\qquad
\omega_{c,rs}=\mu_{rs}\Big[\frac{a_{rc}-a_{rC}}{\mu_r}-\frac{a_{sc}-a_{sC}}{\mu_s}\Big]\quad(c<C,\ r<s),
\]
and its bounds are $\varphi_0+\min/\max\{\omega^\top k:Gk\le g\}$. With $A$ the
sample hull, imposing the rows at every observed $X_i$ is equivalent to imposing
them at the hull's vertices.
\end{proposition}
\begin{reading}
The computable bound is a small LP: one variable per unknown direction and one
constraint per extreme unit, group and outcome. It is an LP for the outer set,
not an LP form of the sharp set, which in general needs infinitely many
conditional constraints.
\end{reading}

\paragraph{Regularity of the value function.} Write
$V(Q,\mu)=\max\{\omega(\mu)^\top k:Gk\le g(Q)\}$. For fixed $\mu$, $V$ is concave
and piecewise linear in $Q$ where the LP is feasible and bounded. For fixed $Q$
it is convex and piecewise linear in $\omega$. Jointly it is in general neither:
$\max\{\mu k:-Q\le k\le Q\}=\mu Q$ is bilinear. $V$ is differentiable where the
optimal basis is unique and non-degenerate, which generalises JKST's
tie-breaking condition. Section~\ref{sec:inference} conditions on the design, so
$\omega$ and $G$ are fixed and only the first case is needed.

\begin{proposition}[Units outside the region]\label{prop:trim}\status{proved}
Let $A$ be a polytope with $P(X\in A)>0$, (S1) everywhere, (S2)--(S3) on $A$,
and A5 on $\supp_AP_X$. Let $\varphi^A_0$ and $\omega^A$ be the quantities of
Proposition~\ref{prop:lp} with $\mu_{rs}$ replaced by $\mu^A_{rs}=\E[NX_rX_s\one\{X\in A\}]$,
$\tilde a(x)_{rc}=a_{rc}x_r/\mu_r$, and
\[
U_{\rm out}=\E\big[N\,h_S(\tilde a(X);X,T)\one\{X\notin A\}\big],\qquad
L_{\rm out}=-\E\big[N\,h_S(-\tilde a(X);X,T)\one\{X\notin A\}\big].
\]
Then $\varphi\in\big[\varphi^A_0+\min_{\ThO}\omega^{A\top}k+L_{\rm out},\
\varphi^A_0+\max_{\ThO}\omega^{A\top}k+U_{\rm out}\big]$, and each endpoint is
attained. For a single cell the outside terms are the Duncan--Davis terms of the
units outside $A$.
\end{proposition}
\begin{reading}
Units outside the region where straight lines are trusted contribute their
model-free range; nothing links them to the model. This is what the package's
\code{range = "trim"} option computes for single cells.
\end{reading}

\subsection{Refutability and undetectable violations}

\begin{proposition}[Non-emptiness and refutation]\label{prop:nonempty}\status{proved}
If $P$ is generated by some $\mathbb Q\in\mathcal M(H_0;A)$ with $P(X\in A)=1$,
then the true $K_0\in\ThI\subseteq\ThO$ and
$B^*_{rc}\in[\varphi_L,\varphi_U]\cap\DD^*_{rc}$, where
\[
\DD^*_{rc}=\Big[\frac{\E[N\max\{0,X_r+T_c-1\}]}{\mu_r},\ \frac{\E[N\min\{X_r,T_c\}]}{\mu_r}\Big].
\]
An empty $\ThO$, or an empty intersection with $\DD^*$, refutes A1--A5 jointly.
\end{proposition}
\begin{reading}
If the model is right, the population bound is never empty. An empty bound is
evidence against the model; in samples it can also be noise, and
Section~\ref{sec:inference} turns it into a test with controlled level.
\end{reading}

\begin{proposition}[What the data can never rule out]\label{prop:undetectable}\status{proved}
Let $\nu:\Delta^{R-1}\to\R^{R\times C}$ be Borel with, at every $x\in\supp P_X$,
(i) $x^\top\nu(x)=0$, (ii) $\nu(x)\one=0$, and (iii) $M_H(x)+\nu(x)\in\Bset$. Then
the structures $\beta=M_H(X)+\nu(X)$ and $\beta'=M_H(X)$ generate the same law of
the data, but
$B^*_{rc}(\beta)-B^*_{rc}(\beta')=\E[NX_r\nu_{rc}(X)]/\mu_r$. Under A5 on $\supp_AP_X$, with $A\supseteq\supp P_X$, the linear
such $\nu$ are exactly $M_{K'}$ with $K'\in\Kspace$.
\end{proposition}
\begin{reading}
There are nonlinear contextual effects that cancel in the aggregate and leave the
data identical, yet change the transition rates. For two rows, $\nu_{1c}=b_cx_1x_2$
and $\nu_{2c}=-b_cx_1^2$ (with $\sum_cb_c=0$) give a quadratic contextual effect
for row 1 that is exactly cancelled in row 2; JKST's Example~6 is the case $C=2$.
The linearity assumption is therefore a genuine leap that no fit statistic can
confirm. Section~\ref{sec:sensitivity} bounds its consequences with a
sensitivity parameter.
\end{reading}

\subsection{Illustration}

In a noise-free $3\times3$ design with 2000 units, rows and columns ordered
(in-party, out-party, non-voter) and an arbitrary linear data-generating process,
the outer bounds on four transition rates are narrower than Duncan--Davis and
contain the truth:

\begin{center}\small
\begin{tabular}{@{}lcccc@{}}
\toprule
Target & Truth & $\ThO$ bound & Duncan--Davis & Width ratio \\
\midrule
Loyalty P$\to$P & 0.848 & [0.396, 0.932] & [0.131, 0.973] & 0.64 \\
Defection P$\to$O & 0.086 & [0.022, 0.361] & [0.000, 0.628] & 0.54 \\
Demobilisation P$\to$N & 0.066 & [0.022, 0.301] & [0.000, 0.532] & 0.52 \\
Mobilisation N$\to$P & 0.169 & [0.048, 0.590] & [0.000, 0.896] & 0.60 \\
\bottomrule
\end{tabular}
\end{center}

The bounds remain wide because six numbers are unidentified. Narrowing them is
the job of the district rungs and of the optional restrictions in the next two
sections. One unknown, for example, is $k_{P,PO}$, half the gap between the
loyalty of in-party voters in a unit otherwise all out-party and the rate at
which out-party voters switch to the in-party in a unit otherwise all in-party;
the data identify the average of these two rates, not the gap.

\paragraph{Relation to Cross and Manski.} In the framework of
\citet{crossmanski2002}, the data-alone region for a unit is the polytope
$S(x,t)$, and their instrument assumption with no further covariates is
Goodman's constancy, the special case of LCE with zero contextual effects,
$H_c=\alpha_c\one^\top$. Their outer region from iterated expectations
corresponds to $\ThO$, and their sharp region to $\ThI$. LCE is weaker than
constancy: rates may vary linearly with composition, so the result is a set
rather than a point. It is also a mean restriction on a superpopulation rather
than an equality of each unit's rates, which is why the bounds concern
aggregates, why each unit stays at its Duncan--Davis polytope, and why the gap
between $B$ and $\Bt$ must be handled in inference.

% Section 4
\section{Districts and the rung ladder}\label{sec:districts}

``District'' means any partition of the units into groups: ridings, counties
within states, schools within boards. Nothing below uses the nesting being
geographic. Let $D\in\{1,\dots,D^\star\}$ with $P(D=d)>0$, $V=(X,N,D)$, and let
$A_d$ be a compact convex set with $P(X\in A_d\mid D=d)=1$; the default is the
district's own composition hull. Write
$\mu^{(d)}_{rs}=\E[NX_rX_s\one\{D=d\}]$ and $\mu^{(d)}_r=\E[NX_r\one\{D=d\}]$, so
that $\mu_r=\sum_d\mu^{(d)}_r$. For a vector $\zeta\in\R^R$, $\zeta\one^\top$ is the
matrix whose row $r$ is constant and equal to $\zeta_r$; since $\one^\top x=1$,
$(\zeta\one^\top)x=\zeta$, so this is an additive shift written in homogeneous form.

The district model replaces (S2) by $\E[\beta\mid V]=M_{H^{(D)}}(X)$ and (S3) by
validity on each $A_d$, with one set of coefficients
$H^{(d)}=(H^{(d)}_1,\dots,H^{(d)}_C)$ per district. Write $Q^{(d)}$ and $K^{(d)}$
for the symmetric and skew parts and $k^{(d)}$ for the free coordinates of
$K^{(d)}$. The rungs are restrictions on $(H^{(1)},\dots,H^{(D^\star)})$, listed in
decreasing order of plausibility:

\begin{center}\small
\begin{tabularx}{\linewidth}{@{}L{0.14\linewidth}Y@{}}
\toprule
Rung & Restriction \\
\midrule
M2 & none: each district has its own linear contextual effects \\
M1($\delta$) & there is $\bar k$ with $\|k^{(d)}-\bar k-c_d\|_\infty\le\delta$ for all $d$, where $c_d$ are the coordinates of $\skw(\zeta_{\cdot,d}\one^\top)$ and $\zeta_{c,d}=\diag(Q^{(d)}_c-Q^{(1)}_c)$ \\
M1 & $H^{(d)}_c=H_c+\zeta_{c,d}\one^\top$ with $\sum_c\zeta_{c,d}=0$ and $\zeta_{c,1}=0$: shared contextual effects, district baselines \\
L4($r$) & M1, and row $r$ has no contextual effect: $H_{c,rs}=H_{c,rr}$ for all $s,c$ \\
L5 & $H^{(d)}_c=\alpha_{c,d}\one^\top$: constancy within each district (Goodman by district) \\
L6 & L5 with $\alpha_{c,d}$ common to all districts (Goodman) \\
M0 & one $H$ for all units on $\conv\bigcup_dA_d$: the pooled model of Section~\ref{sec:rxc}, not nested with the others \\
\bottomrule
\end{tabularx}
\end{center}

The parameter $\delta$ is in probability units. Entries of $K$ are half-gaps
between two hypothetical rates, so $\delta=0.05$ lets the gap between
``$r$-members in pure-$s$ units'' and ``$s$-members in pure-$r$ units'' differ
across districts by up to 10 points.

\subsection{Grouping alone adds nothing}

\begin{proposition}[No model, no gain]\label{prop:D0}\status{proved}
Under (S1) alone, the sharp identified set for $B^*_{rc}$ is the Duncan--Davis
interval $\DD^*_{rc}$ of Proposition~\ref{prop:nonempty}, whatever the law of $D$
or of any other covariate. The joint sharp set for the matrix $B^*$ is convex and
compact but in general not the product of the cell intervals; for a linear
target the sharp interval is
$\big[-\E[N\,h_S(-\tilde a(X);X,T)],\ \E[N\,h_S(\tilde a(X);X,T)]\big]$.
\end{proposition}
\begin{reading}
District dummies cannot narrow the model-free bound, which is already computed
unit by unit. Every gain from districts comes from an assumption that links
them. The accounting identity links the cells within each unit, so the per-cell
extremes cannot in general be attained together, and a combined target can be
strictly narrower than the sum of its cells' intervals.
\end{reading}

\subsection{What each rung identifies}

\begin{lemma}[Identification under the district rungs]\label{lem:D1}\status{proved}
Assume A5 within each district.
(a) Under every rung, $Q^{(d)}$ is identified from $\E[T\mid X,D=d]=X^\top Q^{(d)}X$,
$\sum_cH^{(d)}_c=\one\one^\top$, and each $K^{(d)}\in\Kspace$ is unidentified.
(b) Under M1, $\zeta_{c,d}=\diag(Q^{(d)}_c-Q^{(1)}_c)$ is identified; the
off-diagonal entries satisfy the overidentifying restrictions
$(Q^{(d)}_c-Q^{(1)}_c)_{rs}=\tfrac12(\zeta_{c,d,r}+\zeta_{c,d,s})$ for $r\neq s$;
$K^{(d)}=K+\skw(\zeta_{\cdot,d}\one^\top)$ with $K=K^{(1)}$ common, and the
unidentified dimension equals that of the pooled model.
(c) Under L5, $\alpha_{c,d}=\diag Q^{(d)}_c$, $K^{(d)}=\skw(\alpha_{\cdot,d}\one^\top)$
is identified, and the structure $Q^{(d)}_c=\sym(\alpha_{c,d}\one^\top)$ is testable.
(d) Under M1, A5 within every district can be replaced by a stacked rank
condition: the design made of the pooled quadratic monomials and the district
shifts $\one\{D=d\}X_r$ ($d\ge2$) has full column rank in the population. A
sufficient condition is A5 in one district and $R$ affinely independent
compositions in every other district; $R$ distinct but affinely dependent
compositions are not enough.
\end{lemma}
\begin{reading}
District baselines are identified from each district's own regression. Under
shared effects the curvature of the regression must be the same in every
district: that is the testable part, with $(D^\star-1)(C-1)R(R-1)/2$
restrictions. The shared contextual split $K$ is not testable. Given the common
$Q$, a district's regression is linear in its shift, so a district needs only a
few compositions to fix its baseline, for example two distinct units in
$2\times2$; but with three groups they must not lie on a line. This is
identification, not precision: with about ten units per district the shifts are
identified but noisy.
\end{reading}

\begin{lemma}[The estimand under districts]\label{lem:D2}\status{proved}
Under (S1) and the district version of (S2), with $P(X\in A_D)=1$,
\[
B^*_{rc}=\frac1{\mu_r}\sum_d\Big[\sum_sQ^{(d)}_{c,rs}\mu^{(d)}_{rs}+\sum_{s\ne r}K^{(d)}_{c,rs}\mu^{(d)}_{rs}\Big].
\]
Under M1 its sensitivity to the common $K_{c,rs}$ is $\mu_{rs}/\mu_r$, as in the
pooled model.
\end{lemma}
\begin{reading}
The electorate-wide rate is a weighted sum of district rates. Under shared
effects the unknown enters exactly as in the pooled model, so the district
baselines change the known part of the target, not how much the unknown matters.
\end{reading}

\subsection{Sharp and outer sets for each rung}

\begin{theorem}[Sharpness for affine mean models]\label{thm:affine}\status{proved}
Let a parameter $\theta$ range over a closed set, and let $\mathcal V_A=\{v=(x,n,d):x\in A_d\}$.
Suppose (H1) $v\mapsto m_\theta(v)\in\R^{R\times C}$ is Borel on $\mathcal V_A$;
(H2) $\sup_{v\in\mathcal V_A}\|m_\theta(v)\|<\infty$; (H3) $\theta\mapsto m_\theta(v)$
is affine, and continuous uniformly in $v$ on bounded sets of $\theta$; (H4) the
unit polytopes are not affected by $\theta$. Define the model by (S1),
$\E[\beta\mid V]=m_\theta(V)$ on $\{V\in\mathcal V_A\}$, and $m_\theta(v)\in\Bset$ on
$\mathcal V_A$. Then $\theta$ is in the sharp set iff the validity condition holds
and $\langle u,m_\theta(v)\rangle\le\psi(u;v)$ for $P_V$-almost every $v$ and all $u$.
\end{theorem}
\begin{reading}
The sharpness argument of Theorem~\ref{thm:sharp} carries over to any model in
which the mean is affine in the unknowns, including every district rung.
\end{reading}

\begin{theorem}[The rungs and the ladder]\label{thm:ladder}\status{proved}
Assume A1, A5 within districts and $P(X\in A_D)=1$.
(a) For each rung, the sharp set is the rung's restriction intersected with each
district's own sharp condition (Theorem~\ref{thm:sharp} within district $d$).
(b) The outer set is
$\{\theta\in\text{rung}:(Q^{(d)}_c+K^{(d)}_c(\theta))v\ge0\ \forall d,c,\ v\in\vertx(A_d)\}$.
It is a polyhedron, and every target bound is an LP. For M1 the variables are the
common $k$ and $\Theta^{M1}_O=\bigcap_d\Theta^{(d)}_O$, where $\Theta^{(d)}_O$
imposes district $d$'s rows with its identified shift.
(c) Write $\mathcal B^\rho$ for the set of values of a linear target under rung
$\rho$. For $0\le\delta\le\delta'$,
\[
\mathcal B^{L6}\subseteq\mathcal B^{L5}\subseteq\mathcal B^{L4(r)}\subseteq\mathcal B^{M1}\subseteq\mathcal B^{M1(0)}\subseteq\mathcal B^{M1(\delta)}\subseteq\mathcal B^{M1(\delta')}\subseteq\mathcal B^{M2}\subseteq\DD^*,
\]
for both the sharp and the outer versions, whenever the smaller rung's model is
non-empty.
\end{theorem}
\begin{reading}
Each stronger assumption gives a bound inside the weaker one's, so the ladder can
be reported as a table and readers can see what each assumption buys. The pooled
model M0 is not on this chain. M1(0) keeps the common skew part but drops the
testable symmetric restrictions of Lemma~\ref{lem:D1}(b); it is a sensitivity
benchmark that separates a failure of M1's testable implications from
uncertainty about the unidentified effects, not a substantive rung.
\end{reading}
\begin{sketch}
Each rung is an affine mean model, so Theorem~\ref{thm:affine} applies district
by district. The outer version replaces the sharp condition by non-negativity at
the vertices of each district hull. The nesting holds because each step restricts
the same polyhedral system further; L5 is M1 with constant rows.
\end{sketch}

\subsection{Closed forms and properties}

\begin{proposition}[Shared effects with two rows]\label{prop:D3}\status{proved}
Let $R=2$ and $\vertx(A_d)=\{(t,1-t):t\in\{l_d,u_d\}\}$ with $0<l_d\le u_d<1$, and
let $\lambda_c,\rho_c$ be as in Proposition~\ref{prop:tworows} with $Q=Q^{(1)}$.
Under M1,
\[
\ell_c=\max_d\max_{t\in\{l_d,u_d\}}\Big[\lambda_c(t)-\frac{\zeta_{c,d,1}}{1-t}\Big],\qquad
h_c=\min_d\min_{t\in\{l_d,u_d\}}\Big[\rho_c(t)+\frac{\zeta_{c,d,2}}{t}\Big],
\]
and $\Theta^{M1}_O=\{\kappa:\ell\le\kappa\le h,\ \sum_c\kappa_c=0\}$, projected by
Lemma~\ref{lem:box}. If a district has a pure vertex ($u_d=1$ or $l_d=0$), the
corresponding row at that vertex does not involve $\kappa_c$: it is a testable
restriction and is dropped from the maximum (or minimum), while the other row at
that vertex enters as the formula gives.
\end{proposition}
\begin{reading}
A district raises the lower bound when its baseline for group 1 is low and it has
units that are almost all group 1: there, group 1's rate is nearly observed, and
the shared contextual effect carries that information to every district. One
district with lopsided units and an extreme baseline can tighten the bound for
the whole electorate. The upper bound behaves symmetrically through group 2.
\end{reading}
\begin{sketch}
Under M1, district $d$'s mean for row 1 at vertex $t$ is
$Q_{c,11}t+(Q_{c,12}+\kappa_c)(1-t)+\zeta_{c,d,1}$. Requiring it to be non-negative
and dividing by $1-t>0$ gives a lower limit on $\kappa_c$; row 2 gives an upper
limit after dividing by $t$. The shared $\kappa_c$ must satisfy every district's
limits at every vertex, and the sum-to-zero constraint across columns then gives
a box cut by a hyperplane.
\end{sketch}

\begin{proposition}[When district baselines matter]\label{prop:D4}\status{proved}
(a) If every shift is zero, M1's outer set equals the pooled outer set on
$\conv\bigcup_dA_d$. (b) Under M1 with $P(X\in A_D)=1$, write $Z_d=[\zeta_{c,d,r}]_{rc}$.
Then $\E[\beta\mid X,N]=M_H(X)+\E[Z_D\mid X,N]$, and the pooled model M0 holds
(for some coefficients) iff $\E[Z_D\mid X,N]$ is almost surely a function of $X$
alone that is affine on $\supp P_X$. Independence of $D$ and $(X,N)$ suffices, and
so does conditional independence of $D$ and $N$ given $X$ together with
$\E[Z_D\mid X]$ affine in $X$; independence of $D$ and $X$ alone does not. When M0
holds in this way, its validity condition (S3) is automatic on the population hull
$\conv\supp P_X$ (Lemma~\ref{lem:outer}(c)); on any other pooled region it remains a
separate condition.
\end{proposition}
\begin{reading}
The pooled model is safe when district baselines, given composition and unit
size, vary at most linearly with composition. When strongholds also have
different baselines (a popular local candidate in the strongest districts), or
when district baselines correlate with unit size, the pooled fit mistakes
baseline differences for contextual effects and can miss the truth: an ecological
omitted-variable bias. If two districts have the same distribution of
compositions but one has larger units and a non-zero shift, $\E[Z_D\mid X,N]$
varies with $N$ and the size-weighted pooled target is biased, although the
$X$-only pooled model still holds. Section~\ref{sec:tests} tests the stronger,
design-conditional condition that all shifts vanish.
\end{reading}

\begin{proposition}[District-specific effects]\label{prop:D5}\status{proved}
Under M2 the outer bound is the Minkowski sum of the district bounds,
$\mathcal B^{M2}_O=\sum_d[\varphi^{(d)}_L,\varphi^{(d)}_U]$, where $\varphi^{(d)}$ is
district $d$'s term in Lemma~\ref{lem:D2}. The same holds for the sharp version.
\end{proposition}
\begin{reading}
District-specific lines are safe but rarely much tighter than Duncan--Davis,
because widths add across districts.
\end{reading}

\begin{proposition}[Bounded heterogeneity]\label{prop:D6}\status{proved}
(a) $\delta\mapsto\mathcal B^{M1(\delta)}_O$ is nested. (b) For
$\delta\ge\bar\delta:=\max_d\max_{k\in\Theta^{(d)}_O}\|k-c_d\|_\infty$, which is
finite, it equals $\mathcal B^{M2}_O$. (c) Its endpoints are continuous and
piecewise linear in $\delta$ wherever it is non-empty. (d) At $\delta=0$ it equals
the M1 outer bound computed with the same $Q^{(d)}$ and shifts.
\end{proposition}
\begin{reading}
The dial $\delta$ traces a continuous path between district-specific and shared
effects, so a reader can see how much the shared-effects assumption is doing.
\end{reading}

\begin{proposition}[The most restrictive rungs]\label{prop:D7}\status{proved}
(a) L5 point-identifies the target:
$B^*_{rc}=\sum_d(\mu^{(d)}_r/\mu_r)\,\alpha_{c,d,r}$ with $\alpha_{c,d}=\diag Q^{(d)}_c$,
Goodman's regression within each district. If L5 holds, this point lies in every
larger rung's bound. (b) L4($r$) pins $K_{c,rs}=Q_{c,rr}-Q_{c,rs}$ (reference
district) for all $s\ne r$ and all $c$, which is $(C-1)(R-1)$ of the free
coordinates; for $R=2$ the target is point-identified.
\end{proposition}
\begin{reading}
The constancy rungs give single numbers. Reported beside the bounds, they show
what ``no contextual effects'' implies. If row $r$ does have a contextual effect,
L4($r$) attributes it to the other rows, and the point ($R=2$) or the bound is
misspecified and can exclude the truth, with no warning except an empty set.
With non-voters as row $r$ this is a real risk, because mobilisation often
depends on local campaign intensity. If some district's contextual split differs
from the shared one by more than $\delta$, M1($\delta$) fails in the same way as
M1, but less sharply: the larger $\delta$, the more protection and the less
narrowing.
\end{reading}

\subsection{Worked examples}\label{sec:ladder-examples}

The examples are noise-free, so they concern identification only.

\paragraph{Two groups, two districts.} For in-party loyalty with truth 0.819:
Duncan--Davis gives $[0.453,0.879]$; the pooled model is empty, because the
districts have different baselines; M2 gives $[0.563,0.870]$; and M1 gives
$[0.618,0.866]$. The L5 point is 0.858, which is misspecified, since the true
contextual slopes are 0.10.

\paragraph{Three groups, six districts.} The districts have different
composition ranges and baselines, with a common $H$:

\begin{center}\small
\setlength{\tabcolsep}{3.5pt}
\begin{tabular}{@{}lcccccc@{}}
\toprule
Target & Truth & DD & M2 & M1(.05) & M1 & L5 \\
\midrule
Loyalty P$\to$P & 0.836 & [0.167, 0.953] & [0.391, 0.916] & [0.397, 0.914] & [0.430, 0.880] & 0.870 \\
Defection P$\to$O & 0.095 & [0.000, 0.610] & [0.024, 0.404] & [0.024, 0.390] & [0.042, 0.339] & 0.061 \\
Demobilisation P$\to$N & 0.068 & [0.000, 0.551] & [0.024, 0.350] & [0.024, 0.323] & [0.044, 0.295] & 0.069 \\
Mobilisation N$\to$P & 0.135 & [0.001, 0.685] & [0.045, 0.470] & [0.051, 0.444] & [0.089, 0.413] & 0.122 \\
\bottomrule
\end{tabular}
\end{center}

The ladder shrinks monotonically and contains the truth at every rung. Moving
from M2 to M1 narrows the loyalty bound from width 0.525 to 0.450, and
M1($\delta$) at $\delta=0.10,0.05,0.02$ traces a continuous path between them. The
baseline shifts here are small (standard deviation 0.04); the gain grows with more
heterogeneous baselines and more lopsided units. The L5 point for loyalty is
0.870 against a truth of 0.836: inside the M1 bound, but wrong. The pooled model
is misspecified here, because the baselines vary with composition; it happens to
cover the truth, but Proposition~\ref{prop:D4} gives no guarantee that it will.

% Section 5
\section{Signs, reinforcement and monotonicity}\label{sec:signs}

Monotone ecological inference~\citep{egghm2025} replaces the functional form of
LCE by beliefs about the direction of two associations: whether a group is more
likely than others to choose an outcome within the same unit (the \emph{between}
association), and whether its rate rises with its share across units (the
\emph{within}, or contextual, association). This section states those beliefs in
R$\times$C, shows how they bound the target on their own and on top of a linear
rung, and adds a nonparametric monotone rung.

Two settings are used. In the \emph{finite-population} setting the design is
fixed, unit weights are $p_i=N_i/\sum_jN_j$, $\E_p[f]=\sum_ip_if_i$,
$\mu_r=\E_p[X_r]>0$, and the realised rates satisfy $\beta_i\in S(X_i,T_i)$. The
target is the actual rate $B_{rc}=\sum_iw_{ir}\beta_{irc}$, with
$w_{ir}=p_iX_{ir}/\mu_r$ as before, and the rates are read as real numbers (cell
masses need not be whole numbers of people). In the \emph{superpopulation}
setting, $m(v)=\E[\beta\mid V=v]$, the target is $\Bt$, and a superscript $^*$ on
an association means it is computed with $m$ in place of $\beta$. In this section
$D$ without a unit index denotes a gap between two groups' rates, not a district.

\subsection{The associations and their identities}

Fix a column and drop it from the notation. Write $\Cov_w$ for a covariance under
weights $w$, $w^r_i=w_{ir}$, and $w^{\neg r}_i=p_i(1-X_{ir})/(1-\mu_r)$.

\begin{definition}[Associations]\label{def:assoc}
\begin{enumerate}[label=(\roman*)]
\item Between, pairwise: $\delta_B(r,s)=\E_p[X_rX_s(\beta_r-\beta_s)]$.
\item Between, one-vs-rest: $\delta^r_B=\sum_{s\ne r}\delta_B(r,s)$.
\item Within, pooled (W2): $\Delta=\sum_s\mu_s\Cov_{w^s}(X,\beta_s)\in\R^R$.
\item Within, one-vs-rest (W1): with
$\beta^{\neg r}_i=\sum_{s\ne r}X_{is}\beta_{is}/(1-X_{ir})$, the rate of everyone
outside row $r$,
$\delta^r_W=\mu_r\Cov_{w^r}(\beta_r,X_r)+(1-\mu_r)\Cov_{w^{\neg r}}(\beta^{\neg r},X_r)$.
\end{enumerate}
\end{definition}
\begin{reading}
The between association says that, on average within the same unit, group $r$
is more likely than the others to choose the outcome. It is a weighted average
with weights $p_iX_{ir}(1-X_{ir})$: pure units get no weight and a large mixed
unit can dominate. It is not an ordering in every unit, nor against each rival.
The within association says that rates rise, on average, with group $r$'s share.
In W1 the predictor in both covariances is $X_r$, and the complement's rate is
that of a changing mix of the other groups, so either covariance can be negative
when their sum is positive. A sign on $\Delta_r$ signs one pooled sum of
group-specific covariances, not each group. W1 and W2 can have opposite signs
while every group's rate rises: with groups A, B, C, two units of equal size with
shares $(.2,.7,.1)$ and $(.6,.05,.35)$, and rates $(.40,.80,.10)$ and
$(.42,.82,.12)$, every rate rises by two points, yet $\delta^A_W=-0.02633$ and
$\Delta_A=+0.00110$. Covariances also reflect sorting; a causal reading needs a
separate design. W1 is the one-vs-rest collapse suggested by EGGHM; W2, the
pairwise weighting and reinforcement with those pairings are extensions proved
here.
\end{reading}

Four data-computable quantities appear: Goodman's no-intercept regression
$g=M^{-1}\E_p[XT]$ with $M=\E_p[XX^\top]$; the neighbourhood model
$\NM_r=\E_p[X_rT]/\mu_r$; the one-vs-rest regression slope
$\ER^r=\Cov_p(T,X_r)/\Var_p(X_r)$; and $\NM_{\neg r}=\E_p[(1-X_r)T]/(1-\mu_r)$.
Write $D^r=B_r-B_{\neg r}$ for the gap between group $r$ and the rest, where
$B_{\neg r}=\sum_{s\ne r}\mu_sB_s/(1-\mu_r)$.

\begin{lemma}[Accounting identities]\label{lem:M1}\status{proved}
Assume A1. (a) $\one^\top\Delta=0$ and, if $M$ is invertible, $g-B=M^{-1}\Delta$.
(b) $\NM_r-B_r=-\delta^r_B/\mu_r$. (c) If $\Var_p(X_r)>0$ and $0<\mu_r<1$, the
collapsed population $(X_r,\beta_r,\beta^{\neg r})$ is a valid two-group
population, so
\[
\ER^r-D^r=\frac{\delta^r_W}{\Var_p(X_r)},\qquad(\NM_r-\NM_{\neg r})-D^r=-\frac{\delta^r_B}{\mu_r(1-\mu_r)}.
\]
(d) $B$, $\delta_B(r,s)$, $\delta^r_B$, $\Delta$ and $\delta^r_W$ are affine in
the unit rates, with coefficients that depend only on the design.
\end{lemma}
\begin{reading}
These are accounting identities. The neighbourhood model misses the true gap
exactly by the between association, and the aggregate regression slope misses it
exactly by the within association. So a sign belief about either association is
a one-sided bound relative to a number computed from the data. Signing $\Delta_r$
does not in general sign the bias of row $r$'s Goodman coefficient, because of
the matrix $M^{-1}$.
\end{reading}

\subsection{Sign restrictions on realised tables (N1)}

\begin{definition}[N1 feasible set]
A \emph{sign restriction} is $g_k(\beta)\ge-\varepsilon_k$, where $g_k$ is plus or
minus one of the associations of Definition~\ref{def:assoc} (for any column) and
$\varepsilon_k\ge0$ is a slack (default 0). Optional unit-level restrictions
$\beta_i\in U_i$, with $U_i$ polyhedral, express unit-level orderings or
structural zeros. The \emph{N1 feasible set} is
$F=\{\beta:\beta_i\in S(X_i,T_i)\cap U_i\ \forall i;\ g_k(\beta)\ge-\varepsilon_k\ \forall k\}$.
\emph{Contextual reinforcement} (CR) for a chosen pair $(g_B,g_W)$ of a between
and a within association is $F_{CR}=F_+\cup F_-$ with
$F_\pm=F\cap\{\pm g_B\ge0,\ \pm g_W\ge0\}$: the two associations have the same
sign, whichever it is.
\end{definition}

\begin{proposition}[N1 is sharp for continuous unit tables]\label{prop:N1}\status{proved}
Let $\varphi=\sum a_{rc}B_{rc}$ and read the unit tables as non-negative real cell
masses. (a) The set of values of $\varphi$ attained by continuous unit tables that
reproduce the margins and satisfy the restrictions is $\varphi(F)$; if
$F\ne\emptyset$ it equals $[\min_F\varphi,\max_F\varphi]$. (b) $F=\emptyset$ iff
the restrictions contradict the margins in this continuous problem. (c) Under CR
the set is $\varphi(F_+)\cup\varphi(F_-)$, a union of at most two intervals.
(d) $\varepsilon\mapsto\varphi(F(\varepsilon))$ is nested componentwise, with
continuous piecewise-linear endpoints where $F(\varepsilon)\neq\emptyset$. With
several slacks, the set of slack vectors that preserves a given conclusion is in
general not a product of intervals; its boundary is the breakdown frontier. A
scalar breakdown value needs a stated conclusion and a scalar path
$\varepsilon=t\,v$, $t\ge0$.
\end{proposition}
\begin{reading}
With sign beliefs alone and no regression, the bound is an LP over all unit tables
that fit the margins, and it is exactly right for real-valued tables, conditional
on the observed margins. With whole numbers of ballots it can differ.
\end{reading}

\begin{proposition}[Integer counts, two rows, margins only]\label{prop:integer}\status{proved}
Let $R=C=2$ with no unit-level restrictions, let $N_1$ and $N_0$ be the two
group totals, and $N=N_1+N_0$. The set of gaps $D=B_1-B_2$ attained by integer
tables is the Duncan--Davis interval for $D$ intersected with a lattice of step
$1/N_1+1/N_0=1/(N\mu(1-\mu))$, and both endpoints are attained. Under sign rows
on one column, the integer set is the continuous set of
Proposition~\ref{prop:tworowsign} intersected with that lattice, so the
endpoints differ by less than one step.
\end{proposition}
\begin{reading}
For $N=10^6$ people with $\mu=0.3$ the step is $4.8\times10^{-6}$, so whole
ballots matter only for tiny tables or near-ties. In tiny tables they do matter:
three units of four people with group-1 counts 1, 2, 3 and outcome counts 0, 1, 0
admit the continuous CR set $\{0\}$, but no integer table satisfies CR. With
unit-level restrictions, several columns or $R\ge3$, the integer--continuous gap is
not bounded by this argument; that case is open.
\end{reading}

\begin{proposition}[Two rows]\label{prop:tworowsign}\status{proved}
Let $R=2$, $X=X_1$, $D=B_1-B_2$, $\gamma=\Var_p(X)/(\mu(1-\mu))$, and let
$[D^-,D^+]$ be the Duncan--Davis interval for $D$. Assume $\Var_p(X)>0$, continuous
tables, no unit-level restrictions and no sign rows other than those named. Write
$D_{\ER}$ and $D_{\NM}$ for the one-vs-rest regression slope and neighbourhood
gap.
(a) $\delta_W\ge0\iff D\le D_{\ER}$, and $\delta_B\ge0\iff D\ge D_{\NM}$.
(b) $\{\delta_W\ge0\}$ gives $[D^-,\min\{D^+,D_{\ER}\}]$, empty iff $D_{\ER}<D^-$;
symmetrically, $\{\delta_W\le0\}$ is empty iff $D_{\ER}>D^+$.
(c) $D_{\NM}\in[D^-,D^+]$ always, so a sign on $\delta_B$ is never refuted in this
problem.
(d) $D_{\NM}=\gamma D_{\ER}$ with $\gamma\in(0,1]$. If $\gamma<1$ and
$D_{\ER}\ne0$, exactly one CR branch is non-empty. If $\gamma=1$ (every unit with
positive weight is pure), $D$ is point-identified and both branches are non-empty
with the same gap image $\{D_{\ER}\}$; if $D_{\ER}=0$, both have gap image $\{0\}$.
The CR set is $[D_{\NM},\min\{D^+,D_{\ER}\}]$ if $D_{\ER}\ge0$ and
$[\max\{D^-,D_{\ER}\},D_{\NM}]$ otherwise; under CR,
$\sgn D=\sgn D_{\ER}$ when $D_{\ER}\neq0$.
\end{proposition}
\begin{reading}
In this base problem (two groups, real-valued tables, nothing else imposed) only a
within sign can be contradicted by the margins. A between sign and CR are always
feasible, because ``every group votes like its unit'', $\beta_{i1}=\beta_{i2}=T_i$,
fits them. That non-refutability belongs to this base system only. Under a linear
rung, CR can be rejected jointly with it even when the rung alone fits
(Remark~\ref{rem:LCEexcludesCR}); with integer counts CR can be rejected outright
(Proposition~\ref{prop:integer}); and known cells or unit orderings can remove the
equal-rates witness. In R$\times$C the equal-rates table satisfies every weak
between sign and the CR \emph{union}, but only the branch that matches its within
association.
\end{reading}
\begin{sketch}
By Lemma~\ref{lem:M1}(c) with two rows, $\delta_W=(D_{\ER}-D)\Var_p(X)$ and
$\delta_B=(D-D_{\NM})\mu(1-\mu)$, so each sign is a half-line in $D$, and $D$ ranges
over the whole interval $[D^-,D^+]$. The equal-rates table gives $D=D_{\NM}$. The
identity $D_{\NM}=\gamma D_{\ER}$ with $\gamma\in(0,1]$ shows that the two
same-sign branches of CR overlap only when $\gamma=1$ or $D_{\ER}=0$.
\end{sketch}

\begin{proposition}[Joint restrictions]\label{prop:joint}\status{proved}
Imposing sign restrictions on several columns or rows jointly gives a subset of the
feasible set of any single one of them.
\end{proposition}
\begin{reading}
Joint beliefs are weakly tighter, and often strictly: in $2\times3$ simulations,
signing the within association of two columns jointly narrowed column~1 strictly
in 75 of 139 populations. With two complementary columns, the \emph{same} weak
within sign on both forces both associations to zero and so $D=D_{\ER}$. This is
why the package solves one LP rather than applying $2\times2$ closed forms column
by column.
\end{reading}

\subsection{Signs under a linear rung}

Under a linear rung the model means, for the fixed design, can be written
$m_i=m_{\rm id}(V_i)+L_ik$, where $m_{\rm id}$ is identified, $L_i$ is fixed by the
design, and $k$ stacks the rung's free coordinates: $K$ under M0, the
$K^{(d)}$ under M2, and the common $K$ under M1, with the district shifts in the
identified part.

\begin{proposition}[Sign rows under linear rungs]\label{prop:Krows}\status{proved; (c) for M0 only}
(a) Under every linear rung, $\delta^*_B(r,s)$, $\delta^{r*}_B$, $\Delta^*$ and
$\delta^{r*}_W$ are affine in $k$, so each sign restriction is one linear row in
the LP and CR is a union of two LPs.
(b) If $m_{rc}(x,d)\ge m_{sc}(x,d)$ for every district $d$ and every $x$ in that
district's support, then $\delta^*_B(r,s)\ge0$, pooled over districts.
(c) For $R=2$ under M0, $\delta^*_W=a(w_1+d_1)+b\,w_1$ with $a,b\ge0$, where $w_1$
and $w_1+d_1$ are the two groups' contextual slopes in JKST's notation, so
non-negative slopes imply $\delta^*_W\ge0$. Under the district rungs this
implication fails for the pooled association.
\end{proposition}
\begin{reading}
Under a linear rung, sign beliefs become cheap extra constraints on the unknown
$K$. But a pooled sign under district models is one restriction across all
districts. Positive slopes in every district do not make the pooled within
association positive, because district baselines can fall as composition rises:
by the law of total covariance, the pooled association has a between-district term
with no sign. Example: six equally weighted units with $X=.1,.2,.3$ in district~1
and $.7,.8,.9$ in district~2, both rows' means $.8+.1x$ in district~1 and $.2+.1x$ in
district~2. The slopes are positive, the within-district association is
$+0.00064$, and the pooled one is $-0.04927$.
\end{reading}

\begin{remark}[Realised versus model-implied associations]\label{rem:whichdelta}
N1 uses the realised associations, the linear rows use $\delta^*$. Under (S2),
$\delta-\delta^*=\sum_i\sum_rc_{ir}e_{ir}$ is a design-weighted sum of unit
residuals with mean zero given the design. Under A6 its standard deviation is at
most $\tfrac12\{\sum_i(\sum_r|c_{ir}|)^2\}^{1/2}$; it is small only when no unit
keeps a non-negligible coefficient. A mean condition alone gives no rate.
Intersecting N1 with a linear rung therefore mixes two readings, and the package
reports them separately.
\end{remark}

\subsection{The sign module under the pooled model}\label{sec:msign}

The results of this subsection are proved for the pooled model M0 only. The sign
and CR rows are valid linear rows under the district rungs
(Proposition~\ref{prop:Krows}(a)), but sharpness and inference for them under the
district rungs are open (Section~\ref{sec:open}). The package therefore labels
the module \emph{experimental}.

View $\ThI$ and $\ThO$ in the free coordinates $k$. The target is a model-implied
linear $\varphi(k)=\varphi_0+\omega^\top k$, for example $\Bt_{rc}$ or the gap
$\tilde D=\Bt_1-\Bt_2$. An \emph{admissible restriction} is a set
$C=C(P)\subseteq\R^d$ that is (C1) closed and convex and (C2) depends on $P$ only
through the design and $Q(P)$. Every sign on $\delta^*_B(r,s)$, $\delta^{r*}_B$,
$\Delta^*_r$ or $\delta^{r*}_W$, mean-level ordering on $A$, and the local rows of
Proposition~\ref{prop:sandwich} are admissible and polyhedral. A \emph{union
restriction} is $C_+\cup C_-$ with admissible $C_\pm$; CR for a pair $(g_B,g_W)$
of affine functions of $k$ is $C_\pm=\{\pm g_B\ge0,\pm g_W\ge0\}$.

\begin{theorem}[Sharpness with added restrictions]\label{thm:R3}\status{proved, pooled M0}
Assume A1--A5, $P(X\in A)=1$ and $\affdim(A)=R-1$.
\begin{enumerate}[label=(\alph*)]
\item For admissible $C$, the sharp identified set for $\varphi$ under LCE and $C$
is $[\min_{\ThI\cap C}\varphi,\max_{\ThI\cap C}\varphi]$, or empty.
\item $\ThO\cap C=\bigcup_{P'\in\mathcal P(P)}(\ThI(P')\cap C(P'))$. So the
interval $\varphi(\ThO\cap C)$ is sharp for the regression and the design. It is a
finite LP when $C$ is polyhedral, it equals (a) when $\ThI=\ThO$, and it is weakly
outer otherwise.
\item For a union, the identified set is $\varphi(\Theta\cap C_+)\cup\varphi(\Theta\cap C_-)$
for $\Theta\in\{\ThI,\ThO\}$: at most two intervals. The hull is sharp iff they
intersect, and a common point of $\Theta\cap C_+\cap C_-$ certifies that they do
(sufficient only).
\item For one-vs-rest CR, with $\Var_p(X_r)>0$,
$\gamma_r=\Var_p(X_r)/(\mu_r(1-\mu_r))$, $\ER^{r*}=\Cov_p(X_r,\tau(X))/\Var_p(X_r)$ and
$D^{r*}=\Bt_r-\Bt_{\neg r}$,
\[
C_+\cup C_-=\{k:\ D^{r*}(k)\in\hull\{\gamma_r\ER^{r*},\ \ER^{r*}\}\},
\]
which is two linear rows. For $R=2$ this is the only CR, so in $2\times2$ CR is a
single LP, and under it $\sgn\tilde D=\sgn\ER^*$.
\item Refutation is joint: if $P$ is generated by a model in LCE $\cap\,C$, the
true $k_0\in\ThI\cap C\subseteq\ThO\cap C$; so an empty population LP rejects the
conjunction, not either part alone.
\end{enumerate}
\end{theorem}
\begin{reading}
Adding a sign or ordering belief to the pooled linear rung keeps the bound as
tight as the regression allows: nothing narrower can be learned from the
regression and the design. The full law of the data could rule out more. For the
default one-vs-rest comparison, reinforcement becomes a single band in the gap,
so no union is needed. If the bound is empty, it is the combination that is
contradicted, not the belief on its own. A belief that also has a model-free
meaning can sometimes be tested on its own, and conversely a linear rung can
exclude a belief that is never refutable model-free.
\end{reading}
\begin{sketch}
(a) The identified set is the image of the compact convex set $\ThI\cap C$ under an
affine map; (C2) makes $C$ the same for every law consistent with the data.
(b) Apply the construction of Proposition~\ref{prop:sharpreg} to each point of
$\ThO\cap C$. (d) Apply Lemma~\ref{lem:M1}(c) to the model means:
$\delta^{r*}_W=\Var_p(X_r)(\ER^{r*}-D^{r*})$ and
$\delta^{r*}_B=\mu_r(1-\mu_r)(D^{r*}-\gamma_r\ER^{r*})$; the union of the two
same-sign cases is the hull.
\end{sketch}

\begin{remark}[Full law versus regression]\label{rem:fullreg}
Let $X\in\{1/4,1/2,3/4\}$ with equal probability and equal sizes, and $T$ a fair
binary total at every composition. The full law forces both rows' rates to equal
$T$ in every unit, so $k=0$. The regression alone admits $m_1=.5+k(1-x)$,
$m_0=.5-kx$. For $C=\{k\in[.2,.4]\}$ the full-law set of (a) is empty, while the
regression-only gap set of (b) is $[1/6,1/3]$.
\end{remark}

\begin{remark}[A linear rung can exclude CR although it fits]\label{rem:LCEexcludesCR}
Let $X\in\{0,1/4,1\}$ with equal probability, $A=[0,1]$, $m_0=1-x$,
$m_1=(49/55)x$ and deterministic $T$. The vertex rows force $k=0$, so LCE fits at a
single point, where $\ER=0$ and $\tilde D=-261/1925$. CR requires
$\tilde D\in\{0\}$, so both branches are empty, although LCE alone is not
rejected. In the model-free continuous problem CR is still feasible. Any statement
that a belief ``cannot be refuted'' must name its base system.
\end{remark}

For CR pairings other than W1, part (d) fails: the two sign conditions are on
different affine functions of $k$, so both branches can be non-empty with
disjoint images. An adversarial search found gaps of up to about 31\% of the hull
width. For the pairwise version only $\delta^*_B(r,s)$ is pairwise; $\Delta^*_r$
pools all groups, including those outside the pair.

\begin{proposition}[When does a restriction tighten the bound?]\label{prop:R2}\status{proved}
Let $F$ be the current non-empty compact feasible set, $U=\max_F\varphi$,
$L=\min_F\varphi$, $F^*_U$ and $F^*_L$ the optimal faces, and $N$ an admissible
restriction with $F\cap N\neq\emptyset$ and new endpoints $U_N,L_N$.
\begin{enumerate}[label=(\alph*)]
\item $U_N<U$ iff $F^*_U\cap N=\emptyset$, and $L_N>L$ iff $F^*_L\cap N=\emptyset$.
For polyhedral $N$ each is one feasibility LP, which must carry the target
intercept: $\{Gk\le g,\ \omega^\top k\ge U-\varphi_0,\ k\in N\}$.
\item For one row $a^\top k\le b$: $U_N<U$ iff $\min_{F^*_U}a^\top k>b$, iff every
optimal dual multiplier of the row is positive. A positive multiplier at one
solver optimum does not certify tightening.
\item $N$ is implied ($F\subseteq N$) iff it tightens no linear target.
\item In $2\times2$ under LCE with $\Var_p(X)>0$ and $\gamma<1$, each of the four sign
rows is either implied or tightens every target that is non-constant in the
unknown slope; the closed forms are those of Proposition~\ref{prop:G2}.
\item For a stack of rows, a row's contribution depends on what is already
imposed. Zero gain when added last means zero incremental gain for this target,
not that the row is implied.
\end{enumerate}
\end{proposition}
\begin{reading}
A belief narrows a quantity only if it removes every table that attains that end
of the bound. The package tests this with one feasibility LP per end and never
trusts solver dual values: a row duplicated at an optimal vertex can carry a
positive multiplier without moving anything. A belief can narrow the set of
possible tables without narrowing the quantity of interest; with three groups this
happened in 28\% of simulated cases, against 8\% implied and 64\% tightening.
\end{reading}
\begin{sketch}
Adding a restriction can only lower the maximum. If some maximiser survives, the
maximum is unchanged; if none does, the restricted problem is compact and
non-empty and attains a strictly smaller value. For one row, concavity of the LP
perturbation function ties this to every optimal multiplier being positive. A
hyperplane separating a point of $F\setminus N$ from $N$ gives a target that
tightens, which proves (c).
\end{sketch}

\begin{proposition}[JKST plus signs in $2\times2$]\label{prop:G2}\status{proved}
Under LCE in $2\times2$ with $\Var_p(X)>0$ and $\gamma<1$, let $[L_J,U_J]$ be the
JKST interval for the gap $\tilde D$. Then $\delta^*_W\ge0$ gives
$[L_J,\min(U_J,\ER^*)]$; $\delta^*_W\le0$ gives $[\max(L_J,\ER^*),U_J]$;
$\delta^*_B\ge0$ gives $[\max(L_J,\NM^*),U_J]$; $\delta^*_B\le0$ gives
$[L_J,\min(U_J,\NM^*)]$; and CR gives $[L_J,U_J]\cap\hull\{\NM^*,\ER^*\}$.
Reversed endpoints mean the set is empty.
\end{proposition}
\begin{reading}
Adding a sign belief to JKST cuts its interval at one of two data-computed
numbers. In 2000 random LCE designs with true signs, the within sign tightened the
JKST interval in 83\% of designs (median width ratio 0.70), the between sign in
89\% (0.67), and CR in 80\% of the designs where CR holds (0.69). A sign is
redundant exactly when the rung alone already implies it, that is, when the
interval for that association implied by the rung lies on the believed side of
zero. Under LCE, in $2\times2$ the pair
$(\delta_W,\delta_B)$ lies on a line segment, so re-imposing the rung's own
implications on the associations adds nothing.
\end{reading}

\paragraph{Orderings and local beliefs.} For one column and any $R$, write
$q_{rs}(x)=(Qx)_r-(Qx)_s$, which is identified, and $g_{rs}(x;k)=(Kx)_r-(Kx)_s$,
linear in $k$. The between-group gap at $x$ is $O_{rs}(x)=(Hx)_r-(Hx)_s$, and the
composition-weighted contextual derivative when a member of group $s$ is replaced
by one of group $r$ is $L_{rs}(x)=(x^\top H)_r-(x^\top H)_s$.

\begin{proposition}[Sandwich]\label{prop:sandwich}\status{proved}
(a) $O_{rs}(x)+L_{rs}(x)=2q_{rs}(x)=\partial_{e_r-e_s}\tau(x)$. (b)
$O_{rs}=q_{rs}+g_{rs}$ and $L_{rs}=q_{rs}-g_{rs}$. (c) $O_{rs}\ge0$ and
$L_{rs}\ge0$ at $x$ iff $|g_{rs}(x;k)|\le q_{rs}(x)$, which is possible only if
$q_{rs}(x)\ge0$, a data test.
\end{proposition}
\begin{reading}
An ordering belief (``group $r$ is at least as supportive as group $s$'') and a
local contextual belief restrict the same unknown from opposite sides, so together
they confine it to an interval whose half-width is an identified slope.
\end{reading}

\begin{proposition}[Which directions a pairwise belief can reach]\label{prop:rank}\status{proved}
(a) One pair $(r,s)$ restricts an $(R-1)$-dimensional subspace of the column's
$K$. (b) Pairs forming a star $\{(1,s):s\ne1\}$ reach all $R(R-1)/2$ directions.
(c) Ordering and local rows (Proposition~\ref{prop:sandwich}) on a star bound the
column's $K$ without Duncan--Davis rows.
\end{proposition}
\begin{reading}
With three or more groups, each pairwise comparison can restrict only $R-1$ of
the unknown directions. Enough well-chosen pairs reach all of them; whether any
is actually cut depends on the beliefs and the data. This is the precise sense in
which more groups give more room for beliefs to bind. That every spanning tree of
pairs, not only a star, reaches all directions has been checked numerically for
chains with $R=3,4,5$ but is not proved, and what happens when the support of
compositions is lower-dimensional is open.
\end{reading}

\begin{proposition}[Nesting of beliefs in $2\times2$ under M0]\label{prop:PN}\status{(a), (b) proved}
Let $a=\mu\Var_{w^1}(X)$, $b=(1-\mu)\Var_{w^2}(X)$ and $x^*=a/(a+b)$. (a) Signed
slopes ($w_1\ge0$, $w_1+d_1\ge0$) imply the weighted-local belief $w_1+d_1x\ge0$ on
the hull. (b) $\delta^*_W=(a+b)(w_1+d_1x^*)$, so the weighted-local belief implies
$\delta^*_W\ge0$ whenever $x^*$ lies in the composition hull.
\end{proposition}

\begin{conjecture}[The balance point lies in the hull]\label{conj:C1}
In $2\times2$, $x^*$ always lies in the composition hull.
\end{conjecture}
\begin{reading}
In the pooled $2\times2$ model the weighted-local belief (``replacing a member of
one group by one of the other does not lower the composition-weighted rates, at
every composition'') sits between signed slopes and a sign on the within
association. The second step needs $x^*$ to lie in the observed range.
Conjecture~\ref{conj:C1} says it always does: there was no violation in 2000
random designs and an adversarial search over 50{,}000 discrete designs, but there
is no proof. The chain fails for pooled M1 (Proposition~\ref{prop:Krows}).
\end{reading}

\subsection{Monotone nonparametric contextual effects (N2)}\label{sec:N2}

N2 drops linearity and keeps only direction. It is defined for two rows on the
finite empirical design law: the support of $X$ (the row-1 share) is the set of
distinct observed compositions $x_1<\dots<x_{J_X}$, each with positive mass. Model
N2($\sigma$) assumes (S1), $\E[\beta\mid V]=m(X)$ for a Borel
$m:[0,1]\to\Bset$, and, for designated cells $(r,c)$ with directions
$\sigma_{rc}\in\{+1,-1\}$, that $\sigma_{rc}m_{rc}$ is non-decreasing on the
support. The regression $\tau(x)=\E[T\mid X=x]$ is identified at each support
point.

\begin{lemma}[Means lie in the mean polytope]\label{lem:N21}\status{proved}
$m(x)\in S(x,\tau(x))$ at every support point.
\end{lemma}

Let $\mathcal M_O$ be the set of support-point means with $m(x_j)\in S(x_j,\tau(x_j))$
for all $j$ and
$\sigma_{rc}(m_{rc}(x_{j+1})-m_{rc}(x_j))\ge0$ for all adjacent pairs and designated
cells. Bounds on $\Bt_{rc}=\sum_iw_{ir}m_{rc}(X_i)$ are an LP over $\mathcal M_O$.

\begin{proposition}[The N2 outer set]\label{prop:N2}\status{proved, finite support}
(a) Every $m$ allowed by N2($\sigma$) restricts to a point of $\mathcal M_O$, so the
LP bounds are an outer set for $\Bt$; they are not claimed sharp.
(b) Every point of $\mathcal M_O$ extends to a monotone Borel $m$ on $[0,1]$, so on
the complete finite support nothing is lost by restricting to support points. On a
coarser grid this fails, and the grid LP is only a relaxation.
(c) If $\tau$ is quadratic, M0 with signed slopes lies inside N2, for the outer sets.
(d) If $m_{1c}$ and $m_{2c}$ are both non-decreasing, $\delta^*_W\ge0$ for column $c$.
(e) For $R\ge3$ a direction ``in $x_1$'' needs a chart saying which share adjusts,
and (d) has no general analogue.
\end{proposition}
\begin{reading}
N2 assumes only ``no reversals'': each designated group's mean rate moves in one
direction as its share grows, in any shape. It is weaker than straight lines with
signed slopes, so its bound contains theirs, and the two-row shape restriction is
stronger than a pooled covariance sign, which permits reversals. If a designated
mean reverses direction somewhere (a U-shape or a tipping point), the order rows
exclude the true mean function, and the bound can be narrow, non-empty and wrong;
an empty set is the only warning, and plotting the fitted regression for
reversals is a partial check. N2 is not in the package.
\end{reading}

\begin{conjecture}[The sharp N2 set]\label{conj:N2}
The sharp N2 set is the set of $m$ with $m(X)\in\E[S(X,T)\mid V]$ $P_V$-almost
surely, plus the shape restrictions, where the expectation is the conditional
selection expectation given $V=(X,N)$.
\end{conjecture}
\begin{reading}
When unit size carries information, conditioning on composition alone admits mean
rates that the full model excludes. Fix $X=1/2$ and two equally likely sizes; at
the first size $T$ is 0 or 1 with equal probability, at the second $T=1/2$ surely.
At the first size both rows' rates must equal $T$, so the model forces
$m(1/2)=(.5,.5)$; the polytope LP, which conditions on $X$ only, admits
$(.75,.25)$. Under the selection-expectation reading the characterisation is
sufficient on a finite support: for each design atom take a selection whose
conditional mean is the common candidate and glue the selections across atoms.
Necessity in general is not written, so the statement stays a conjecture. A
procedure that uses only $\E[T\mid X]$ does not attain the set when size is
informative.
\end{reading}

\begin{proposition}[Local identity]\label{prop:N23}\status{proved, as a calculus identity}
Let $m_{1c},m_{2c}$ be differentiable at $x$, and let
$\tau_{m,c}(z)=z\,m_{1c}(z)+(1-z)m_{2c}(z)$ be the structural aggregate,
$D_c(x)=m_{1c}(x)-m_{2c}(x)$ and $\delta_{W,c}(x)=x\,m'_{1c}(x)+(1-x)m'_{2c}(x)$.
(a) $\tau'_{m,c}(x)=D_c(x)+\delta_{W,c}(x)$. (b) $\delta_{W,c}(x)\ge0\iff
D_c(x)\le\tau'_{m,c}(x)$. (c) Local CR, $D_c(x)\delta_{W,c}(x)\ge0$, holds iff
$D_c(x)$ lies between 0 and $\tau'_{m,c}(x)$. (d) If both rows are non-decreasing on
a neighbourhood of $x$, the local sign at $x$ is non-negative; monotonicity on the
finite support is not enough, and the local sign everywhere neither implies nor is
implied by the sign of $\delta^*_W$.
\end{proposition}
\begin{reading}
EGGHM's local results are statements about smooth functions, but on a finite
support the data never identify a derivative: adding
$\prod_j(x-x_j)$ to both rows leaves the law unchanged and moves the derivative at
every node. So local sign beliefs cannot be imposed from aggregates without an
identified regression family, and they are not an observable rung here.
\end{reading}

\begin{lemma}[Person-weighted bins]\label{lem:bins}\status{proved}
Let $b=1,\dots,K$ be ordered non-overlapping bins of compositions with
person-weighted masses $P(b)>0$ and within-bin expectations $\E_b$. Put
$\bar x_b=\E_bX$, $\bar m_{1b}=\E_b[Xm_1(X)]/\bar x_b$,
$\bar m_{2b}=\E_b[(1-X)m_2(X)]/(1-\bar x_b)$ and $\bar\tau_b=\E_b[\tau(X)]$, the
model-implied bin mean. Under N2($\sigma$): (i) $\bar m_b\in S(\bar x_b,\bar\tau_b)$;
(ii) the targets are linear in the bin rates with known weights; (iii) monotonicity
is preserved across the ordered bins. The binned LP is an outer relaxation of the
N2 bounds.
\end{lemma}
\begin{reading}
With many compositions, group them into bins, but average rates by weighting the
people in each group, not the units: the ordinary bin average breaks the
accounting identity (it gives 0.25 against the correct 0.3125 in a two-point
example). In samples the observed person-weighted bin mean of $T$ estimates
$\bar\tau_b$ with error, which inference must account for.
\end{reading}

\paragraph{What inference for N2 needs.} The regression at the support points
must be estimated. A weighted polynomial of \emph{prespecified} degree $q$ makes
the constraints linear in its coefficients, and the confidence intervals of
Section~\ref{sec:inference} then apply conditional on the degree. The degree must
not be chosen by cross-validation and then treated as fixed: in a balanced
two-level design with true slope zero, choosing between an intercept-only and a
linear fit by leave-one-out error and then using the nominal 95\% slope interval
gives conditional coverage of about 0.68.

% Section 6
\section{Inference}\label{sec:inference}

\subsection{Sampling, estimators and regimes}

\paragraph{Where the randomness comes from.} \citet{aaiw2023} distinguish
sampling-based uncertainty (the units are a sample from a larger population) from
design-based uncertainty (some assignment is random). EI has neither: every unit
of the electorate is observed, and nothing is assigned. In a purely design-based
view nothing is random, $B$ is a fixed unknown, and the only valid bound is
Duncan--Davis. LCE, however, is a statement about a conditional expectation,
which has content only if unit rates are random given the unit's
characteristics. The uncertainty here is therefore \emph{model-based, conditional
on the design}, as in the fixed-regressor framework of \citet{aiz2014}:
\begin{itemize}
\item[(I1)] Given the design and the district effects, the unit rates $\beta_i$ are
independent with $\E[\beta_i\mid\text{design}]=m(V_i)$ (Assumption A6).
\item[(I2)] Inference is conditional on $\{(N_i,X_i,D_i)\}$, so the vertices of
the sample hulls are fixed.
\end{itemize}
District shocks are fixed features of this electorate. Under M1 and M2 they are
district effects, part of the mean. Under M0 a district shock is not a variance
problem but a misspecification: conditional on the design, M0's mean model holds
for the observed units only if the district shifts are zero
(Proposition~\ref{prop:tests}(c)), and clustered standard errors would not
repair it. No clustered variance is used. Independence given the district
effects is a substantive assumption, to be argued in each application: within-district
shocks can correlate residuals with no district mean shift, which violates (I1)
without violating the mean model. Residual correlation within districts and
between neighbours is a diagnostic, not a test that can confirm (I1). Under
dependence, the variance and residual-allowance statements below are not
covered.

\paragraph{Estimators.} Fit each of the $C-1$ free columns by OLS on the
homogeneous quadratic monomials $X_rX_s$ ($r\le s$). For M1 add the district shifts
$\one\{D=d\}X_r$ ($d\ge2$), which imposes the overidentifying restrictions of
Lemma~\ref{lem:D1}(b); for M2 fit one regression per district. Column $C$ follows
by adding up. Write $\hat\beta$ for the stacked coefficients, $\beta_0$ for their
true value and $\hat\Sigma$ for an HC3 covariance estimate. Each LP constraint is
$s_j(\beta)+a_j^\top k\ge0$, with $s_j$ linear in $\beta$ and $a_j$ fixed by the
design (Proposition~\ref{prop:lp}, Theorem~\ref{thm:ladder}(b)). The target is
$\Bt(k)=\varphi_0(\beta)+\omega^\top k$, with $\varphi_0$ linear in $\beta$ and
$\omega$ computed from the sample weights $w_{ir}$, so population moments are
never estimated. The \emph{plug-in bound} is
$\varphi_0(\hat\beta)+[\min,\max]\{\omega^\top k:\ s(\hat\beta)+Ak\ge0\}$. It is an
estimate of the identified set, not a confidence interval.

\begin{table}[t]
\centering\small
\begin{tabularx}{\linewidth}{@{}L{0.12\linewidth}L{0.15\linewidth}Y L{0.25\linewidth}@{}}
\toprule
Model & Regime & Result & Status \\
\midrule
any rung & any, any $n$ & Theorem~\ref{thm:finite} (finite-sample interval and emptiness test) & theorem; uninformative in width so far \\
M0 & any & Propositions~\ref{prop:consistency}, \ref{prop:triangular}, Theorems~\ref{thm:limit}, \ref{thm:ci} & theorem (asymptotic) \\
M1, M2 & R1 ($D^\star$ fixed) & the same & theorem (asymptotic) \\
M1, M2 & R3 & the same, with a high-dimensional Gaussian approximation~\citep{cck2017} under a rate condition & conjecture \\
M1 & R2 & coverage of the CLR interval & simulation finding \\
any & any & the CLR interval with simulated candidates & Conjecture~\ref{conj:clr} \\
\bottomrule
\end{tabularx}
\caption{Which inference result holds in which regime.}\label{tab:regimes}
\end{table}

Under M0 the number of constraints grows with the number of hull vertices, but
the regression has fixed dimension. Every constraint term is a linear functional
of one asymptotically normal vector, the events $\{\max_j(\cdot)\le c\}$ are convex
sets, and the multivariate central limit theorem holds uniformly over convex sets
\citep[Theorem~17.5]{bhattacharyarao1976}. So the growth in constraints does not
matter. Under M1 and M2 in R2 or R3, $\dim\beta$ grows with $D^\star$ and this
argument no longer applies.

\subsection{The structure of the bound and its plug-in estimator}

\begin{lemma}[An endpoint is a minimum of linear functions]\label{lem:minlin}\status{proved, fixed constraint system}
If the LP is feasible with bounded value at $\beta$, then
\[
\varphi_U(\beta)=\varphi_0(\beta)+\min_{\lambda\in\Lambda}\lambda^\top s(\beta),\qquad
\Lambda=\vertx\{\lambda\ge0:A^\top\lambda=-\omega\},
\]
a finite set. So $\varphi_U$ is concave and piecewise linear in $\beta$; likewise
$\varphi_L(\beta)=\varphi_0(\beta)-\min_{\lambda\in\Lambda'}\lambda^\top s(\beta)$ with
$\Lambda'=\vertx\{\lambda\ge0:A^\top\lambda=\omega\}$.
\end{lemma}
\begin{reading}
Each end of the bound is the tightest of finitely many candidate limits, each
linear in the regression coefficients; write $\theta_\lambda(\beta)=g_\lambda^\top\beta$
for the candidate with gradient $g_\lambda$. JKST's case has four candidates and
assumes one wins clearly. In R$\times$C with districts there are many candidates,
and near-ties are the norm. The lemma concerns a fixed constraint matrix and
objective, which conditioning on the design provides.
\end{reading}

\begin{proposition}[Consistency and rate]\label{prop:consistency}\status{proved}
Assume (I1)--(I2), A1--A5 for the rung used, $\sqrt n(\hat\beta-\beta_0)=O_p(1)$,
and Slater's condition: some $\tilde k$ and $\varepsilon>0$ satisfy
$s_j(\beta_0)+a_j^\top\tilde k\ge\varepsilon$ for every constraint defined by the
population hulls. Then (a) $\hat\varphi_U-\varphi_U(\beta_0;\hat A)=O_p(n^{-1/2})$,
and likewise for $\varphi_L$; (b) $\varphi_U(\beta_0;\hat A)\to\varphi_U(\beta_0;A)$
when $\hat A\to A$ in Hausdorff distance along a nested sequence, or along any
sequence whose feasible sets eventually lie in a common bounded set; (c) since
$\hat A_d\subseteq A_d$, $\varphi_U(\beta_0;\hat A)\ge\varphi_U(\beta_0;A)$ and
$\varphi_L(\beta_0;\hat A)\le\varphi_L(\beta_0;A)$.
\end{proposition}
\begin{reading}
The plug-in bound converges at the usual rate to the bound for the sample hull,
which is weakly wider than the population bound; estimating the hull only ever
widens it. The hull error is one-sided and tends to zero, but no rate is shown, so
the limit theory below concerns the sample-hull endpoint. Coverage of the true
parameter never depends on the hull, because the true parameter satisfies the
constraints at every observed composition. Slater's condition fails when every
feasible $k$ makes some implied rate exactly 0 at some vertex, a structural zero,
which is not exotic in EI; the rate result is then open, but the coverage results
below do not use Slater. Bounds do not shrink to a point: their width converges to
the population width.
\end{reading}

\begin{theorem}[Limit of the plug-in endpoint, fixed constraint system]\label{thm:limit}\status{cited structure}
Under the assumptions of Proposition~\ref{prop:consistency} and
$\sqrt n(\hat\beta-\beta_0)\Rightarrow\mathcal Z\sim N(0,\Sigma)$, suppose the
constraint system is fixed along the sequence, so that $A$, $\omega$, $\Lambda$ and
the gradients $g_\lambda$ do not depend on $n$. Let $\Lambda^*$ be the set of optimal
dual vertices at $\beta_0$. Then
\[
\sqrt n\big(\hat\varphi_U-\varphi_U(\beta_0)\big)\Rightarrow\min_{\lambda\in\Lambda^*}g_\lambda^\top\mathcal Z .
\]
(i) If $\Lambda^*$ is a singleton, the limit is normal. (ii) If two active
gradients differ in a direction in which the estimator varies,
$(g_\lambda-g_{\lambda'})^\top\Sigma(g_\lambda-g_{\lambda'})>0$, the limit has
strictly negative mean: the plug-in upper endpoint is biased downward, the lower
endpoint upward, and the plug-in bound is too narrow. If all active gradients
agree on the support of $\mathcal Z$, the limit is normal. (iii) In case (ii) the
standard bootstrap is inconsistent.
\end{theorem}
\begin{reading}
When several candidate limits are nearly tied, the plug-in end picks the one that
looks tightest by chance, so the estimated bound is biased inward. This is why the
plug-in is not a confidence interval. The proof applies the delta method for
directionally differentiable maps \citep{shapiro1991,fangsantos2019} and the
bootstrap characterisation of \citet{fangsantos2019}. With two tied candidates and
unit noise, the mean plug-in error is $-1/\sqrt\pi\approx-0.564$.
\end{reading}

\begin{remark}[Use $C-1$ columns]\label{rem:Cminus1}
The column shares sum to one, so the coefficients of the $C$ column regressions
sum to a known constant and the joint covariance of all $C$ columns is singular.
Stating $\hat\beta$, $\Sigma$ and the gradients for the first $C-1$ columns, and
recovering column $C$ by adding up, lets $\Sigma$ be non-singular. When it is
positive definite, the condition in (ii) reduces to $g_\lambda\ne g_{\lambda'}$.
\end{remark}

\begin{proposition}[Changing constraint systems]\label{prop:triangular}\status{proved}
Let the design, and hence the constraint system, change with $n$, with dual vertex
sets $\Lambda_n$, $|\Lambda_n|\le\bar L$, bounded gradients $g_{\lambda,n}$, and
localised slacks $c_{\lambda,n}=\sqrt n(\theta_{\lambda,n}(\beta_0)-\varphi_{U,n}(\beta_0))\ge0$.
Suppose that, after relabelling, $g_{\lambda,n}\to g_\lambda$ and
$c_{\lambda,n}\to c_\lambda\in[0,\infty]$. Then
\[
\sqrt n\big(\hat\varphi_{U,n}-\varphi_{U,n}(\beta_0)\big)\Rightarrow\min_{\lambda:\,c_\lambda<\infty}\big(g_\lambda^\top\mathcal Z+c_\lambda\big).
\]
The limit of Theorem~\ref{thm:limit} obtains when every $c_\lambda\in\{0,\infty\}$.
\end{proposition}
\begin{reading}
As the sample hull grows, candidates can move in and out of contention. A
candidate that is inactive only by a gap of order $n^{-1/2}$ still shapes the
limit. For example, $f_n(\theta)=\min\{\theta,-\theta+n^{-1/2}\}$ at $\theta_0=0$ has a
uniquely active first piece for every $n$, yet
$\sqrt n(f_n(\hat\theta)-f_n(\theta_0))\Rightarrow\min\{Z,1-Z\}$, which is not normal.
The result describes the sample-hull endpoint and says nothing about the
population-hull error. Uniform limit theory without convergent localised slacks is
open.
\end{reading}
\begin{sketch}
By Lemma~\ref{lem:minlin} the scaled error equals, exactly,
$\min_\lambda\big(g_{\lambda,n}^\top\sqrt n(\hat\beta-\beta_0)+c_{\lambda,n}\big)$.
Candidates with $c_{\lambda,n}\to\infty$ drop out with probability tending to one,
and the continuous mapping theorem gives the limit.
\end{sketch}

\subsection{Confidence intervals}\label{sec:ci}

Fix levels $\alpha_U,\alpha_L,\alpha_E\in(0,1)$. Let $\hat\sigma_\lambda$ be the
standard error of $\hat\theta_\lambda=\theta_\lambda(\hat\beta)$ computed from
$\hat\Sigma$, $\hat\sigma_s$ the vector of standard errors of the constraint terms
$s_j(\hat\beta)$, and $\hat\sigma_{\varphi_0}$ that of $\varphi_0(\hat\beta)$. For a cell
$(r,c)$ write $w_i=w_{ir}$, $S_1=\tfrac12(\sum_iw_i^2)^{1/2}$ (Popoviciu's bound
on the standard deviation of the residual term) and $z_{\alpha_E}=\Phi^{-1}(1-\alpha_E/2)$,
where $\Phi$ is the standard normal distribution function.

\begin{theorem}[Confidence intervals]\label{thm:ci}
\begin{enumerate}[label=(\alph*)]
\item \status{proved} \emph{Conservative interval.} Let $\hat c$ be the
$(1-\alpha_U-\alpha_L)$-quantile of the maximum of the standardised Gaussian errors of
all constraint terms (one-sided) and of the target constant (two-sided). Put
$\hat\Theta^+=\{k:s(\hat\beta)+\hat c\hat\sigma_s+Ak\ge0\}$ and
\[
\CI=\Big[\varphi_0(\hat\beta)-\hat c\hat\sigma_{\varphi_0}+\min_{\hat\Theta^+}\omega^\top k-z_{\alpha_E}S_1,\ \
\varphi_0(\hat\beta)+\hat c\hat\sigma_{\varphi_0}+\max_{\hat\Theta^+}\omega^\top k+z_{\alpha_E}S_1\Big].
\]
If $\max_iw_i^2/\sum_iw_i^2\to0$, then
$\liminf_nP(B_{rc}\in\CI)\ge1-\alpha_U-\alpha_L-\alpha_E$. Without the $S_1$ terms
the interval covers $\Bt_{rc}$ with probability at least $1-\alpha_U-\alpha_L$.
\item[(b1)] \status{cited} \emph{CLR with the full candidate set.} Let $\Lambda$ be
the full set of dual vertices of Lemma~\ref{lem:minlin}, finite for a fixed
design, and assume Condition~P of \citet{clr2013}: $\Sigma$ positive definite
(using $C-1$ columns and a full-rank design) and $g_\lambda\ne0$ for every
$\lambda$. For the upper end, select
$\hat V=\{\lambda:\hat\theta_\lambda\le\min_\mu(\hat\theta_\mu+\hat c_p\hat\sigma_\mu)+2\hat c_p\hat\sigma_\lambda\}$,
where $\hat c_p$ is the $1-0.1/\log n$ quantile of $\max_\lambda Z_\lambda$ and $Z$ is
the standardised Gaussian vector with correlations from $\hat\Sigma$; set
$\hat U=\min_{\lambda\in\hat V}(\hat\theta_\lambda+\hat c\hat\sigma_\lambda)$ with $\hat c$
the $(1-\alpha_U)$-quantile of $\max_{\lambda\in\hat V}Z_\lambda$; treat the lower end
alike; and add $z_{\alpha_E}S_1$ on each side. The coverage bound of (a) holds
\citep[Theorem~4]{clr2013}.
\item[(b2)] \status{proved} \emph{Any candidate set with a simultaneous critical
value.} Let $\hat\Lambda$ be any data-dependent set of dual-feasible points, and
$\hat U=\min_{\lambda\in\hat\Lambda}(\hat\theta_\lambda+\hat c\hat\sigma_\lambda)$, with
$\hat U=+\infty$ if $\hat\Lambda$ is empty. Then
$\liminf_nP(\varphi_U(\beta_0)\le\hat U)\ge1-\alpha_U$ in either of two cases:
(i) $\hat c$ is the $(1-\alpha_U)$-quantile of $\max_{\lambda\in\bar\Lambda}Z_\lambda$
for a fixed finite set $\bar\Lambda\supseteq\hat\Lambda$, with $\sigma_\lambda>0$ on
$\bar\Lambda$; (ii) \emph{projection}: $\hat c=(\chi^2_{q,1-\alpha})^{1/2}$ and
$\hat\sigma_\lambda=(g_\lambda^\top\hat\Sigma g_\lambda/n)^{1/2}$, where either
$\Sigma$ is positive definite and $q=\dim\beta$, or there is an exact reduced
full-rank parameterisation $\beta=L\eta+b$, $\hat\beta=L\hat\eta+b$, with $q=\dim\eta$,
$\Sigma_\eta$ positive definite and $\hat\Sigma=L\hat\Sigma_\eta L^\top$. In case
(ii) one event covers both ends, for any candidate sets, at total level $\alpha$.
\end{enumerate}
\end{theorem}
\begin{reading}
The intervals widen each candidate limit by its own uncertainty before taking the
tightest, so noisy candidates (for example, districts with few units) rarely set
the limit. Three versions carry proofs. The conservative interval widens all
constraints at once. The projection interval uses one ellipsoid for the whole
coefficient vector; it is valid for any collection of candidates, however chosen,
but it is wide when there are many coefficients, because its critical value grows
like $\sqrt q$. The fixed-superset version is sharper but needs a fixed superset.
For the actual rate $B$, an extra allowance covers each unit's local luck. JKST's
interval $\CI_x$ is version (b) with a single candidate per side and no selection
step, with the two error sources combined by adding standard deviations.

An exact reduction means documented identities, such as the adding-up identity
that removes column $C$, or coordinates fixed by the design. A rank read off
$\hat\Sigma$ with a numerical cutoff is not one. A singular limit does not control
errors of smaller order in its null directions: with
$\hat\beta-\beta_0=(Z/\sqrt n,-1/n)$, the second coordinate has standard error zero
and the bound $\hat\beta_2+0$ misses $\beta_{0,2}$ in every sample. So the projection
interval is stated with $q=\dim\beta$ unless an exact reduction is documented.
Dropping candidates that load on numerically null directions is then only a
safeguard: it can only widen the interval.
\end{reading}
\begin{sketch}
(a) On one simultaneous event, every true constraint term is at most its inflated
estimate and the target constant is within its allowance, so the true $k_0$ lies in
$\hat\Theta^+$ and $\Bt$ lies in the inflated LP range. The residual term $E$ has
conditional variance at most $S_1^2$ by Popoviciu's inequality; along
subsequences, Chebyshev's inequality handles the case where its actual standard
deviation is negligible relative to $S_1$, and the Lindeberg central limit theorem
the rest. A union bound finishes. (b2) By weak duality every dual-feasible
candidate is an upper bound for the true endpoint, so it suffices that all
candidates are simultaneously covered: by the fixed-superset maximum in (i), and by
Cauchy--Schwarz in the metric of $\hat\Sigma_\eta$ in (ii).
\end{sketch}

\begin{conjecture}[CLR with simulated candidates; the package default]\label{conj:clr}
Collect $\hat\Lambda$ as the optimal dual vertices of the LP at $\hat\beta$ and at
$n_{\rm draw}$ Gaussian draws around it, and compute both $\hat c_p$ and $\hat c$
over $\hat\Lambda$ in the rule of Theorem~\ref{thm:ci}(b1). Under the conditions
of (b1), the resulting interval has asymptotic coverage at least $1-\alpha_U$ for
the upper end (and likewise for the lower end) as $n_{\rm draw}\to\infty$.
\end{conjecture}

The procedure is an \emph{algorithm}: its validity is not proved. What is proved
is a reduction.

\begin{lemma}[Reduction to the full-candidate rule]\label{lem:reduction}\status{proved}
Suppose $\hat c_p$ is computed over the full set $\Lambda$, and $\hat\Lambda$
contains every vertex selected by the rule of Theorem~\ref{thm:ci}(b1) over
$\Lambda$ and a minimiser of $\hat\theta_\mu+\hat c_p\hat\sigma_\mu$ over $\Lambda$.
Then the threshold, the selected set, $\hat c$ and $\hat U$ coincide with those of
(b1), so its coverage carries over.
\end{lemma}
\begin{reading}
Three steps remain open: recovery of vertices that are optimal only for rare
Gaussian directions (and degenerate optimal faces); the fact that computing
$\hat c_p$ over a subset of $\Lambda$ can only lower it, an error in the
anti-conservative direction; and the growth of $|\Lambda|$ with the number of
districts, outside CLR's fixed-index setting. Until these are settled, the CLR
interval as implemented is an algorithm with simulation support
(Section~\ref{sec:evidence}), and the intervals of Theorem~\ref{thm:ci}(a) and
(b2) carry the formal guarantee.
\end{reading}

\paragraph{A variance-adaptive residual allowance.} Popoviciu's bound uses the
worst-case variance $\tfrac14$ for every unit. The Bhatia--Davis inequality
\citep{bhatiadavis2000} gives $\Var(e_{irc}\mid V_i)\le m_i(1-m_i)$ with
$m_i=m_{rc}(V_i)$. Put $\pi_i=w_i^2/\sum_jw_j^2$ and $\bar m=\sum_i\pi_im_i$, the
$\pi$-weighted model-implied rate, which is itself a linear target. By concavity of
$m(1-m)$,
\[
s_n^2:=\Var(E_{rc}\mid\text{design})\le\sum_iw_i^2m_i(1-m_i)\le\Big(\sum_iw_i^2\Big)\bar m(1-\bar m)=:\bar s^2 .
\]
Let $\CI_v$ be a confidence interval for $\bar m$ at level $1-\alpha_v$, and
$\hat S_{BD}=\big(\sum_iw_i^2\max_{m\in\CI_v\cap[0,1]}m(1-m)\big)^{1/2}$.

\begin{proposition}[Variance-adaptive allowance]\label{prop:bd}\status{proved}
Suppose $\CI_v$ is an interval of Theorem~\ref{thm:ci}(a) or a projection interval
of Theorem~\ref{thm:ci}(b2)(ii), computed on the rung's own constraint system
(without added sign rows), and $\max_iw_i^2/\bar s^2\to0$. Replacing $z_{\alpha_E}S_1$
by $z_{\alpha_E}\hat S_{BD}$ in the conservative interval, or in the projection
interval for $B$, gives
$\liminf_nP(B_{rc}\in\CI)\ge1-\alpha_{\rm est}-\alpha_E-\alpha_v$, where
$\alpha_{\rm est}$ is the level spent on the estimation part.
\end{proposition}
\begin{reading}
The allowance for local luck can be made smaller when the plausible rates are far
from one half, provided the input interval for $\bar m$ comes from a proved
method. The condition on the weights holds, for instance, when $\bar m$ stays
bounded away from 0 and 1. If $\CI_v$ fails, the allowance is not conservative:
with true $\bar m=.5$, a wrong interval inside $[0,.1]$ gives a variance bound of at
most $.09$ against the true $.25$. If $\CI_v$ is instead the CLR interval, the
same bound holds conditional on $\CI_v$ covering $\bar m$, and the result inherits
the status of Conjecture~\ref{conj:clr}. An optional, stronger assumption (the
unit-level deviations of different groups are not negatively correlated) would
make the residual variance estimable from the outcome shares and remove most of
the allowance; it is not in the default.
\end{reading}
\begin{sketch}
On the event that $\CI_v$ covers $\bar m$, $\hat S_{BD}\ge\bar s\ge s_n$. So a
miss of the residual allowance requires either that $\CI_v$ fails or that
$|E|>z_{\alpha_E}\bar s$, whose limiting probability is at most $\alpha_E$ by the
Chebyshev--Lindeberg argument with $\bar s$ in place of $S_1$. A union bound with
the estimation event finishes.
\end{sketch}

\paragraph{Boundary cases.} When the true $K$ is on the boundary of the identified
set, the side of the interval that touches it carries the full error probability,
and coverage is close to nominal. When the true $K$ is interior, coverage tends to
one: this is the distinction between covering the parameter and covering the
set~\citep{imbensmanski2004}. When the bound's width is near zero (L4 in
$2\times2$, L5), one-sided intervals per endpoint under-cover unless the critical
value adapts~\citep{stoye2009}. The two-sided split used here is uniformly valid
at some cost in width.

\subsection{Tests of the model}\label{sec:tests}

\begin{proposition}[Tests]\label{prop:tests}
(a) \status{proved} \emph{Emptiness test.} Rejecting the rung when $\hat\Theta^+$ is
empty has asymptotic level at most $\alpha_U+\alpha_L$.
(b) \status{proved for fixed dimension} \emph{Shared against district-specific
effects.} A Wald test of the M1 regression against M2 tests the
$(D^\star-1)(C-1)R(R-1)/2$ overidentifying restrictions of Lemma~\ref{lem:D1}(b);
its $\chi^2$ calibration holds for fixed dimension (regime R1).
(c) \status{proved} \emph{Pooling.} Conditional on the design, suppose each
district's observed compositions contain $R$ affinely independent points. Under
M1, M0's mean model holds for the observed units iff all district shifts vanish.
Pooling is therefore tested by a Wald test of $\zeta=0$ in the M1 regression, with
$(D^\star-1)(C-1)R$ restrictions.
\end{proposition}
\begin{reading}
An empty confidence set is a real rejection at the stated level. The curvature
test checks the testable part of shared effects. The pooling test says whether the
pooled model is wrong for this design; a rejection does not mean that its
intervals fail (Section~\ref{sec:evidence}), and the test has no power against
within-district dependence without mean shifts. Proposition~\ref{prop:D4}(b) gives a
weaker condition for pooling, but it averages over district labels and so
describes a superpopulation of re-drawn districts, which is not the target here;
conditional on the observed labels only $\zeta=0$ makes M0 correct.
\end{reading}
\begin{sketch}
(a) On the simultaneous event of Theorem~\ref{thm:ci}(a), the true $k_0$ lies in
$\hat\Theta^+$, so emptiness requires that event to fail. (c) If M0 held for the
observed units with coefficients $H'$, then in each district
$(G_c-\zeta_{c,d}\one^\top)X_i=0$ for all its units, with $G=H'-H$. The district's
$R$ affinely independent compositions lie in $\{\one^\top x=1\}$ and are therefore
linearly independent, so $G_c=\zeta_{c,d}\one^\top$ in every district. With
$\zeta_{c,1}=0$, $G_c=0$ and every shift vanishes.
\end{sketch}

\paragraph{Calibration by a restricted wild bootstrap.} In moderate dimension the
HC3 Wald statistic compared with $\chi^2$ over-rejects. The package calibrates
both tests by a restricted wild bootstrap: fit the null model (M0 for the pooling
test, M1 for the curvature test); draw $Y^*=\hat Y_0+\hat e_0v$ with Rademacher
$v_i$, one per unit, shared across the $C-1$ equations; refit the larger model and
recompute the same Wald statistic; and report
$p=(1+\#\{W^*\ge W\})/(B+1)$ with $B=399$ by default. The $\chi^2$ and F p-values are
reported as well. This calibration is supported by simulation
(Section~\ref{sec:evidence}); its validity as the number of restrictions grows is
not proved.

\subsection{A finite-sample interval valid in every regime}

Fix a rung and condition on the design. Let $Z$ be the rung's regression design,
with full column rank, and $\hat\beta_c=(Z^\top Z)^{-1}Z^\top T_{\cdot c}$ for
$c<C$. Every linear functional $f(\beta)=f_0+\sum_{c<C}b_c^\top\beta_c$ satisfies
$f(\hat\beta)=f_0+\sum_i\sum_{c<C}w_{ic}(f)T_{ic}$ with
$w_{ic}(f)=b_c^\top(Z^\top Z)^{-1}z_i$. Put $w_{iC}(f)=0$,
$\rho_i(f)=\max_cw_{ic}(f)-\min_cw_{ic}(f)$, and the Hoeffding radius
$h(f;\delta)=\big(\tfrac12\log(1/\delta)\sum_i\rho_i(f)^2\big)^{1/2}$. With $J$
constraint rows and $\delta=\alpha_s/(J+2)$, put $h_j=h(s_j;\delta)$,
$h_0=h(\varphi_0;\delta)$, $t_E=\big(\tfrac12\log(2/\alpha_E)\sum_iw_{ir}^2\big)^{1/2}$,
$\hat\Theta^{\rm fin}=\{k:s_j(\hat\beta)+h_j+a_j^\top k\ge0\ \forall j\}$, and
\[
\CI^{\rm fin}=\Big[\varphi_0(\hat\beta)-h_0+\min_{\hat\Theta^{\rm fin}}\omega^\top k-t_E,\ \
\varphi_0(\hat\beta)+h_0+\max_{\hat\Theta^{\rm fin}}\omega^\top k+t_E\Big].
\]

\begin{theorem}[Finite-sample interval and emptiness test]\label{thm:finite}\status{proved}
Assume (I1), that the rung's mean model holds and that $Z$ has full column rank.
Then for every $n$ and in every regime: (a)
$P(B_{rc}\in\CI^{\rm fin}\mid\text{design})\ge1-\alpha_s-\alpha_E$; (b) without $t_E$
the interval covers $\Bt_{rc}$ with probability at least $1-\alpha_s$; (c)
rejecting the rung when $\hat\Theta^{\rm fin}$ is empty is a test of exact level at
most $\alpha_s$; (d) intersecting with Duncan--Davis keeps (a).
\end{theorem}
\begin{reading}
This interval needs no Gaussian approximation, no fixed dimension, no Slater
condition and no regime, so it is the only formal guarantee for shared effects with
many small districts. In every design tried it is as wide as the logical range
$[0,1]$. Two costs combine: Hoeffding's inequality uses the worst-case variance
for each outcome share (the true variance was about 0.01, against the worst case
$\tfrac14$), and the
constraint-by-constraint union bound inflates each radius, which is then divided
by the distance from the vertex to a pure composition. An oracle Gaussian union
bound with the true variances still gave widths of 0.76--0.87, so most of the cost
is the union bound, which the conservative asymptotic interval shares. The theorem
settles validity in regime R2, but not usefulness; it is not in the package.
\end{reading}
\begin{sketch}
Every constraint term and the target constant are linear functionals of the OLS
fit, and each is a sum over units of independent terms whose range is known from
the design, because the outcome shares lie in the simplex. OLS is conditionally
unbiased under the mean model, so Hoeffding's inequality~\citep{hoeffding1963}
gives a finite-sample radius for each. A union bound over the $J+2$ one-sided
events puts the true $k_0$ in the widened set, and the residual term gets its own
Hoeffding radius.
\end{sketch}

\subsection{Many districts with few units}

\begin{proposition}[Dropping and selecting districts]\label{prop:select}\status{proved}
(a) For any subset $\mathcal D'$ of districts, the M1 outer set computed from the
constraints of $\mathcal D'$ contains the full M1 outer set. (b) If $\mathcal D'$ is
chosen from the design only, any procedure with (finite-sample or asymptotic)
coverage for a fixed set of districts keeps that coverage conditionally on the
selected $\mathcal D'$; in particular Theorem~\ref{thm:finite} holds. The asymptotic
results transfer only if their own regime conditions hold for the retained
districts.
\end{proposition}
\begin{reading}
With many small districts, ignoring some districts' constraints can only widen the
bound, so it never costs validity of the identified set. If the retained districts
are chosen from the design alone (for example, those with at least 30 units or
enough spread in composition), not from the outcomes, any interval that is valid
for a fixed set of districts stays valid. Selection does not by itself move the
problem into a regime where the asymptotic theory applies: retained districts with
few units are still in R2. Other practical options are M1($\delta$) with $\delta>0$,
which bounds how far one district can pull the shared parameter, coarser groups,
and targeting $\Bt$ to drop the residual term. A hierarchical model for the district
shifts would change the model and is not pursued.
\end{reading}

\subsection{Sensitivity to nonlinearity}\label{sec:sensitivity}

\begin{proposition}[Bounded nonlinearity]\label{prop:S1}\status{proved; an outer bound}
Replace (S2) by $\E[\beta\mid V]=M_H(X)+\nu(X)$ with $\nu$ in the undetectable class
of Proposition~\ref{prop:undetectable} ($x^\top\nu(x)=0$, $\nu(x)\one=0$) and
$\|\nu(x)\|_\infty\le\eta$. Then $Q$ is identified as before, and every value of
$\Bt_{rc}$ consistent with the data lies in
\[
\Big[\min_{\ThO(\eta)}\varphi-\eta,\ \max_{\ThO(\eta)}\varphi+\eta\Big],\qquad
\ThO(\eta)=\{k:s+Ak\ge-\eta\one\}.
\]
For a general target $\sum a_{rc}B_{rc}$, widen by $\eta\sum|a_{rc}|$ instead.
\end{proposition}
\begin{reading}
If rates may depart from straight lines, invisibly, by up to $\eta$, the bound
widens in a controlled way. The bound is valid, but it is not claimed to be the
sharp set of the relaxed model. The \emph{breakdown value} $\eta^*$ of a
conclusion, such as ``loyalty exceeds 0.30'', is the largest departure under which
it survives; the bounds are nested in $\eta$, so it can be found by bisection. In
the $3\times3$ illustration of Section~\ref{sec:rxc}, $\eta^*(\text{loyalty}>0.30)=0.060$:
the conclusion survives any undetectable nonlinearity of up to six percentage
points. Applying the confidence-interval version of the bound accounts for
sampling error. The dial $\delta$ of Proposition~\ref{prop:D6} works the same way
for heterogeneity across districts.
\end{reading}

\subsection{Inference with sign rows, unions and near-empty sets}\label{sec:R10}

The sign rows of Section~\ref{sec:msign} enter the same machinery. A row
$s_j(\beta)+a_j^\top k\ge0$ with $s_j$ linear in $\beta$ and $a_j$ fixed by the design
is \emph{linear in $\beta$}; every sign, ordering and local row has this form,
because the associations, the targets and $\ER^*$ are linear in the model means
with design weights, and $Q$ is linear in $\beta$. A \emph{parallel disjunction} is
a row $\max_{l\le L}s_l(\beta)+a^\top k\ge0$, whose pieces share the normal $a$. A
\emph{general union} is $C_+\cup C_-$, each branch a finite set of rows linear in
$\beta$. $\CI(\mathcal R)$ denotes an interval of Theorem~\ref{thm:ci} computed with the
base rows plus the rows $\mathcal R$. For version (a), $\CI=\emptyset$ iff
$\hat\Theta^+=\emptyset$; for version (b), iff the computed ends cross.

\begin{theorem}[Inference with added rows]\label{thm:R10}\status{proved for pooled M0, by interval version}
Assume the conditions of Theorem~\ref{thm:ci} for the augmented row system, the
pooled model M0, and that the restrictions hold at the truth.
\begin{enumerate}[label=(\alph*)]
\item \emph{Added rows.} For finitely many rows linear in $\beta$,
Theorem~\ref{thm:ci} and Proposition~\ref{prop:tests}(a) hold for the augmented
system, version by version (Table~\ref{tab:R10}). For an actual target that is a
contrast, such as $B_1-B_0$, the residual allowance must cover the contrast:
$\tfrac12\{\sum_i(w_{i1}+w_{i0})^2\}^{1/2}$ is a valid standard-deviation envelope,
and so is the sum of the per-cell envelopes.
\item \emph{Parallel disjunction} (conservative version only). Inflate each piece
separately and then take the maximum,
$\max_l\big(s_l(\hat\beta)+\hat c\hat\sigma_l\big)+a^\top k\ge0$, with all pieces in the
critical-value maximum. This covers at the nominal level. One-vs-rest and
$2\times2$ CR are of this form (the \emph{slab}), and no branch is selected from
$\widehat{\ER}$.
\item \emph{General union.} $\CI(C_+)\cup\CI(C_-)$, each branch at the full level
$\alpha$, covers the true target with asymptotic probability at least $1-\alpha$. No
Bonferroni split is needed. This is point coverage; it gives no simultaneous
coverage of both branch images.
\item \emph{Near-empty sets.} Under the model,
$\limsup_nP(\CI=\emptyset)\le\alpha_U+\alpha_L$, including point-identified cases,
and conditioning on non-emptiness can only raise coverage. If the restrictions are
false, no coverage statement holds.
\item \emph{Plug-in brackets are not valid.} Imposing the estimated bracket
$\hull\{\gamma_r\widehat{\ER},\widehat{\ER}\}$ without inflating it can under-cover,
and more data do not cure it: in a repeated three-point design,
$\limsup_nP(\tilde D_0\in\CI_{\rm naive})\le\Phi(\tfrac7{12}z_{.995})=0.9335<0.95$.
\end{enumerate}
\end{theorem}

\begin{table}[t]
\centering\small
\begin{tabularx}{\linewidth}{@{}L{0.32\linewidth}Y L{0.17\linewidth}@{}}
\toprule
Interval version & What Theorem~\ref{thm:R10} claims & Status \\
\midrule
(a) conservative & (a)--(d), including the slab (b) & proved \\
(b1) CLR over all dual vertices & (a), (c), (d), under CLR's conditions for the full augmented system & cited \\
(b2)(i) fixed superset & the branch-union and crossing arguments & proved \\
(b2)(ii) projection & both ends, any candidates; $\Sigma$ positive definite or an exact reduction & proved \\
CLR with simulated candidates (package default for CR) & the inclusion and crossing steps are deterministic; their probabilities are conjectural & algorithm (Conjecture~\ref{conj:clr}) \\
\bottomrule
\end{tabularx}
\caption{Status of Theorem~\ref{thm:R10} by interval version.}\label{tab:R10}
\end{table}

\begin{reading}
Sign beliefs enter the same interval machinery as the rung's own constraints.
Reinforcement uses a band whose two pieces are each widened by their own
uncertainty before the larger is taken; this is valid without knowing which branch
is true. The tempting shortcut, plugging the estimated neighbourhood and regression
numbers into the bracket, is not a confidence interval. Intersecting the fitted
bracket with a finished JKST interval is worse: its limiting coverage of a truth on
the bracket's edge is at most one half. When the belief is false, the interval stays
narrow and can be reproducibly wrong (Section~\ref{sec:evidence}); a narrow,
non-empty interval is no evidence that the restrictions are true. Conditioning on
other data events, or selecting a belief because it gives an attractive result, is
not covered.
\end{reading}
\begin{sketch}
(a) The added rows are linear in the regression coefficients, so they join the same
simultaneous Gaussian event, on which the true $k_0$ satisfies every inflated row.
(b) The piece that holds at the truth is dominated on the same event, so no branch
is ever chosen from the data. (c) The true branch is a fixed population fact; the
interval for that branch alone has the full level and lies inside the union.
(d) An empty interval requires the coverage event to fail. (e) In the three-point
design, with the basis $((1-x)^2,2x(1-x),x^2)$, the target constant's standard error
is exactly $\tfrac7{12}$ of the regression slope's, and the critical value is at
most $z_{.995}$ by the union bound over the eight base rows and the two target
tails.
\end{sketch}

The statements are pointwise in the data-generating process. Uniform validity near
point identification is open. Statements of the type of
Proposition~\ref{prop:consistency} and Theorem~\ref{thm:limit} need Slater's
condition, which rows that bind at the truth (such as the edge of the CR band) can
break; the conservative version does not.

\subsection{A practical recipe}

\begin{enumerate}
\item Fit the rung's regression on $C-1$ columns with HC3 covariance, and compute
the sample hulls.
\item Build the LP and report the plug-in bound, labelled as an estimate of the
identified set.
\item Report the CLR interval for $\Bt$ and for $B$, labelled as algorithmic, and
beside it the projection interval as the interval with an asymptotic guarantee;
for $B$, take $\bar m$ from the projection interval on the rung alone, or use
Popoviciu's bound. Intersect both with Duncan--Davis, which always contains $B$.
\item With many small districts, say that the CLR interval's coverage rests on
simulation. Theorem~\ref{thm:finite} gives a finite-sample guarantee, but so far
only at the width of the logical range.
\item Report the emptiness test and, for M1, the curvature and pooling tests with
wild-bootstrap p-values. A rejection of pooling says that M0 is wrong for this
design, not that its intervals fail.
\item Report sensitivity curves in $\eta$ and $\delta$, with the breakdown value of
the headline conclusion.
\item With few units per district, select districts by design, or prefer
M1($\delta$) or coarser groups, and say so.
\end{enumerate}

% Section 7
\section{The R package}\label{sec:package}

The methods are implemented in the R package \pkg{eiPartialID}, version 0.3.1,
which extends the $2\times2$ package of JKST. The original functions
(\code{generateBounds()}, \code{evaluateBounds()}) are unchanged; with a range
offset of $10^{-5}$ and HC1 variances, the new code reproduces them on JKST's
census example to $10^{-10}$. The LPs are solved with \pkg{lpSolve} and the
robust covariances come from \pkg{sandwich}.

\subsection{The main function}

\code{eiBounds()} computes bounds and confidence intervals under one rung. Its
main arguments and defaults are:
\begin{itemize}
\item the data, as a formula or as \code{x}, \code{t} and unit sizes \code{n}, with
an optional \code{group} (the districts);
\item \code{model}: \code{"M0"} (default), \code{"M1"}, \code{"M1delta"} (with
\code{delta}), \code{"M2"}, \code{"L4"} (with \code{row\_constant}), \code{"L5"},
\code{"L6"};
\item \code{targets}: linear targets as weight matrices (by default every cell rate);
\item \code{range}: \code{"hull"} (default; the group-wise sample hull,
Proposition~\ref{prop:consistency}(c)), \code{"trim"} (dropped units enter
through Duncan--Davis, Proposition~\ref{prop:trim}) or \code{"extrapolate"} (pure
units or a user matrix, with a warning that validity there is assumed,
Proposition~\ref{prop:simplex});
\item \code{weights}: \code{"none"} or \code{"n"} for the regression;
\code{vcov}: \code{"HC3"} (default), \code{"HC1"}, \code{"HC0"};
\item \code{ci}: \code{"clr"} (default; Conjecture~\ref{conj:clr}, with
\code{ndraw = 100} simulated draws), \code{"conservative"}
(Theorem~\ref{thm:ci}(a)) or \code{"none"}; \code{alpha = 0.05};
\item \code{fpc}: the residual allowance for the actual rate,
\code{"bhatia-davis"} (default; Proposition~\ref{prop:bd}) or \code{"popoviciu"}
($S_1$);
\item \code{restrictions}: optional rows (\code{zeros}, \code{order}, \code{sign},
\code{slope}, \code{linear}, \code{external}), and the sign module's \code{assoc}
(signs on $\delta_B$, one-vs-rest or pairwise, and on $\delta_W$, W1 or W2) and
\code{cr} (contextual reinforcement);
\item \code{union}: \code{"pieces"} (print every piece of a union) or \code{"hull"};
\item \code{eta}: the sensitivity parameter of Proposition~\ref{prop:S1};
\code{jkst = TRUE} adds JKST's $\CI_x$ with its revert-to-Duncan--Davis rule;
\code{tests = TRUE} computes the tests of Proposition~\ref{prop:tests}.
\end{itemize}

\paragraph{What a fit reports.} For each target, \code{eiBounds()} returns:
\code{estimate}, the plug-in estimate of the outer identified set for the
model-implied rate (not a confidence interval); \code{ci\_model}, the CLR (or
conservative) interval for $\Bt$; \code{ci\_actual}, the same for $B$ with the
residual allowance; \code{proj\_model} and \code{proj\_actual}, the projection
intervals of Theorem~\ref{thm:ci}(b2)(ii) computed from the same candidates, with
$q=\dim\beta$ and, for $B$, the Bhatia--Davis allowance taken from a projection
interval for $\bar m$ on the rung alone (Proposition~\ref{prop:bd}); and the
Duncan--Davis interval. Intervals for $B$ are intersected with Duncan--Davis. The
rank of $\hat\Sigma$ is printed as a diagnostic; candidates whose gradients load on
a numerically null direction of $\hat\Sigma$ are dropped as a labelled safeguard,
which can only widen the projection interval. A CLR candidate with zero estimated
variance is treated as exact only if its gradient is zero.

With optional restrictions, the rung alone is also fitted and printed beside the
restricted result; with several restrictions it is labelled as their joint effect.
The \code{tests} component holds the emptiness test, the curvature test (M1 against
M2) and the pooling test, each with $\chi^2$, F and restricted-wild-bootstrap
p-values ($B=399$). An empty estimate is reported as a sample diagnostic, not a
test; the emptiness test is the test.

\paragraph{What \code{print()} says about each fit.} Every fit prints a status line
that names the result covering its intervals:
\begin{itemize}
\item \emph{Pooled M0:} the projection intervals carry an asymptotic guarantee
(Theorem~\ref{thm:ci}(b2)(ii)), or the conservative interval does
(Theorem~\ref{thm:ci}(a)) when \code{ci = "conservative"}; with sign or CR rows the
line adds that the guarantee is Theorem~\ref{thm:R10} for the pooled model.
\item \emph{Grouped rungs without sign rows:} an asymptotic guarantee only with few,
large groups (regime R1); with many small groups, a simulation finding.
\item \emph{Grouped rungs with sign or CR rows:} no coverage theorem.
\end{itemize}
CLR lines are labelled ``algorithmic, coverage supported by simulation, not
proved''; for $B$ the label adds that the Bhatia--Davis bound also uses a CLR
interval; with CR, the CLR interval is the union of the CLR branch intervals.
Every sign or CR restriction prints a line marking the module as experimental,
with its proof status. Grouped fits with sign rows print that each sign is one
pooled restriction across groups and that the bound is an outer set whose
sharpness under grouped rungs is not proved; pooled fits may label the estimate
sharp for the regression (Theorem~\ref{thm:R3}(b)). Notes flag near-boundary fits
(Slater slack below 0.005), few units per group (median below 30, where CLR
coverage rests on simulation), unbounded estimates (a hull that is not
full-dimensional), a Bhatia--Davis input interval for $\bar m$ inside $[0,.05]$
or $[.95,1]$ (where the condition of Proposition~\ref{prop:bd} is doubtful and
\code{fpc = "popoviciu"} is recommended), and joint rejections, naming the rejected
combination and saying whether the rung alone is rejected.

\subsection{Other functions}

\begin{itemize}
\item \code{eiRedundancy()} reports what each restriction buys, using the face test
of Proposition~\ref{prop:R2} at the point estimate (never solver duals): the
verdict and the width gain when the restriction is added first (to the rung alone)
and last (to the full stack); the associations already settled by the rung
alone, with their intervals; the CR branches and the connectivity certificate of
Theorem~\ref{thm:R3}(c); EGGHM's descriptive quantities ($\gamma$, $\ER$, $\NM$),
with a warning never to intersect the bracket with a confidence interval; and the
weight concentration of each signed association (the largest unit weight and the
effective number of units).
\item \code{ladder()} fits the rungs from the most to the least plausible
(Theorem~\ref{thm:ladder}(c)).
\item \code{sensitivity(dial = "eta")} and \code{breakdown()} implement
Proposition~\ref{prop:S1}; \code{sensitivity(dial = "slack")} relaxes the sign
beliefs by a slack in probability units, with a breakdown value along a stated
scalar path (Proposition~\ref{prop:N1}(d)).
\item \code{eiModules()} lists the restriction types with their proof status, and
\code{eiBenchmark()} computes coverage on data with known truth.
\end{itemize}

\paragraph{Not implemented.} Sign rows on realised tables (N1), the monotone rung
N2, local contextual rows, the series bound that would tighten the outer set
towards $\ThX$, the finite-sample interval of Theorem~\ref{thm:finite} (because it is
uninformative), covariates and integer tables.

% Section 8
\section{Numerical evidence and an application}\label{sec:evidence}

The results below are evidence, not proof. They check that the theory behaves as
stated in designs where the truth is known, and they show where the proved
intervals are informative. All intervals are nominal 95\%.

\subsection{Verification}

Every closed form, inclusion and counterexample in Sections~\ref{sec:rxc}--\ref{sec:inference}
has a numerical or symbolic check against brute force or an LP solver. Examples:
the homogeneous form reproduces JKST's bounds to $10^{-16}$; the two-row closed
form agrees with the LP in 300 random cases, and its district version
(Proposition~\ref{prop:D3}) in 126 cases plus 195 with pure vertices; the face test
of Proposition~\ref{prop:R2} agrees with brute force in 800 of 800 cases; and every
counterexample quoted in the readings (the singleton hull, the null-direction
projection failure, the changing-system limit $\min\{Z,1-Z\}$, the
district-reversal example, the three-point bracket design) is reproduced exactly.
Two parts are not machine-checked: the measure-theoretic steps of
Theorems~\ref{thm:sharp} and~\ref{thm:affine}, which rest on cited theorems, and the
remark on randomising over integer vertices, which is only a sketch and is not used.

\subsection{Monte Carlo design}

Unit rates are Beta (for $2\times2$) or Dirichlet (for $3\times3$) draws whose
conditional means are exactly linear in composition given the district. So M1
always holds, and M0 holds only without district shifts. Each design has 1000
replications. Coverage is the share of replications whose interval contains the
truth; an empty interval (an empty inflated set, or crossing ends) counts as not
covering. A coverage of 1 means every replication covered, and the exact one-sided
95\% lower confidence bound is then $0.05^{1/1000}=0.997$. Monte Carlo standard
errors are in parentheses. Widths are medians after clipping to $[0,1]$. The
projection interval for $B$ takes its Bhatia--Davis allowance from a projection
interval for $\bar m$ on the rung alone at level 0.005, with 0.04 on the
estimation part (Proposition~\ref{prop:bd}); the CLR interval for $B$ uses the
Bhatia--Davis allowance with a CLR input. Designs:
\begin{itemize}
\item A: $2\times2$, M1, 10 districts of 93 units (regime R1) and 93 districts of 30;
\item C: $2\times2$ pooled, with an implied rate close to 0 at one vertex, so
Slater's condition holds, but only just;
\item E: $2\times2$ pooled, compositions spanning $[0.02,0.98]$, with group $b$'s
mean rate exactly 0 at the right vertex and group $w$'s exactly 0 at the left, so
the identified set of the model-implied rate is a single point and Slater fails;
\item B: $3\times3$ pooled, four targets.
\end{itemize}

\begin{table}[t]
\centering\footnotesize
\setlength{\tabcolsep}{3pt}
\begin{tabular}{@{}lrc ccc ccc@{}}
\toprule
& & & \multicolumn{3}{c}{Model-implied rate $\Bt$} & \multicolumn{3}{c}{Actual rate $B$} \\
\cmidrule(lr){4-6}\cmidrule(lr){7-9}
Design & $q$ & DD & Conservative & CLR & Projection & Conservative & CLR & Projection \\
\midrule
A, $10\times93$ & 21 & 0.574 & 1 / 0.596 & 1 / 0.432 & 1 / 0.554 & 1 / 0.696 & 1 / 0.538 & 1 / 0.662 \\
A, $93\times30$ & 187 & 0.567 & 1 / 0.633 & 1 / 0.414 & 1 / 0.932 & 1 / 0.694 & 1 / 0.477 & 1 / 0.964 \\
C, boundary & 3 & 0.325 & 0.996 / 0.041 & 0.994 / 0.039 & 0.999 / 0.041 & 1 / 0.087 & 0.999 / 0.057 & 1 / 0.059 \\
E, Slater fails & 3 & 0.271 & 0.990 / 0.179 & \textbf{0.944} / 0.143 & 0.995 / 0.194 & 1 / 0.249 & 0.998 / 0.207 & 1 / 0.251 \\
B, loyalty & 12 & 0.835 & 1 / 0.592 & 1 / 0.557 & 1 / 0.581 & 1 / 0.675 & 1 / 0.659 & 1 / 0.679 \\
B, defection & 12 & 0.635 & 1 / 0.399 & 1 / 0.369 & 1 / 0.395 & 1 / 0.447 & 1 / 0.431 & 1 / 0.445 \\
B, demobilisation & 12 & 0.533 & 1 / 0.333 & 1 / 0.304 & 1 / 0.329 & 1 / 0.381 & 1 / 0.363 & 1 / 0.378 \\
B, mobilisation & 12 & 0.884 & 1 / 0.619 & 1 / 0.559 & 1 / 0.595 & 1 / 0.683 & 1 / 0.646 & 1 / 0.661 \\
\bottomrule
\end{tabular}
\caption{Monte Carlo, model holds: coverage / median width (1000 replications). $q=\dim\beta$; DD is the median Duncan--Davis width. Design E, CLR for $\Bt$: Monte Carlo standard error 0.007.}\label{tab:mc1}
\end{table}

\begin{table}[t]
\centering\footnotesize
\setlength{\tabcolsep}{4pt}
\begin{tabular}{@{}rrccccccc@{}}
\toprule
& & & Plug-in & \multicolumn{3}{c}{Model-implied rate $\Bt$} & \multicolumn{2}{c}{Actual rate $B$} \\
\cmidrule(lr){5-7}\cmidrule(lr){8-9}
Districts & $q$ & DD & cover / empty & CLR & Projection & Theorem~\ref{thm:finite} & CLR & Conservative \\
\midrule
30 & 61 & 0.571 & 0.305 / 0.241 & 0.999 / 0.557 & 1 / 0.967 & 1 / 1.000 & 1 / 0.746 & 1 / 1.000 \\
93 & 187 & 0.566 & 0.042 / 0.463 & 0.995 / 0.467 & 1 / 1.000 & 1 / 1.000 & 0.998 / 0.575 & 1 / 0.930 \\
200 & 401 & 0.566 & 0.005 / 0.717 & 0.996 / 0.416 & 1 / 1.000 & 1 / 1.000 & 1 / 0.491 & 1 / 0.890 \\
\bottomrule
\end{tabular}
\caption{Monte Carlo, many small districts (regime R2): M1, $2\times2$, 10 units per district. Coverage / median width. The projection interval and Theorem~\ref{thm:finite} for $B$ have coverage 1 and width 1.000 in every row.}\label{tab:mc2}
\end{table}

\subsection{Monte Carlo results}

Tables~\ref{tab:mc1} and~\ref{tab:mc2} report the designs in which the model holds.

\begin{enumerate}
\item \emph{No proved interval under-covers where its theorem applies.} The
projection interval covers at least 0.995 (design E, standard error 0.002) and the
conservative interval at least 0.990 (design E, 0.003). Theorem~\ref{thm:finite}
covered in every replication of every design where it was computed, including
regime R2 and designs A. In the R2 rows ($93\times30$ and 10 units per district),
$q$ grows with the number of districts, so the coverage of the projection and
conservative intervals there is itself a simulation finding.
\item \emph{CLR (a simulation finding) covers at least 0.994 in every design where
Slater's condition holds}, including R2 with 200 districts. At the Slater failure
(design E), CLR for $\Bt$ covers 0.944 (standard error 0.007), within one standard
error of nominal. This is the informative case: the identified set is a point, the
interval has no slack, and coverage falls from near one to the nominal level. The
coverage of 1 elsewhere reflects truths in the interior of the identified set.
\item \emph{The plug-in fails as Theorem~\ref{thm:limit} and
Proposition~\ref{prop:triangular} predict.} In R2 its coverage falls to 0.305, 0.042
and 0.005 at 30, 93 and 200 districts, and it is empty 24\%, 46\% and 72\% of the
time. In design E it covers 0.267 and is empty half the time: with a point
identified set, any sampling noise empties it.
\item \emph{The projection interval is informative only when $q$ is small.} Its
critical value grows like $\sqrt q$: 2.8 at $q=3$, 4.6 at $q=12$, 14.8 at $q=187$ and
21.2 at $q=401$, against about 2--3 for CLR. In pooled designs it is only 0.01--0.03
wider than CLR; with 10 districts it is 0.554 against 0.432; with 93 districts it
is 0.93, and $[0,1]$ at 10 units per district. In grouped designs with many districts,
CLR is the only informative interval, and its support there is simulation only.
\item \emph{Theorem~\ref{thm:finite} is valid but uninformative}: clipped width
1.000 in every M1 design, including $10\times93$, where CLR is 0.43 and
Duncan--Davis 0.57.
\item \emph{In R2, CLR narrows as districts are added}: for $\Bt$, 0.557, 0.467 and
0.416 at 30, 93 and 200 districts. With 200 districts of 10 units it is narrower
than Duncan--Davis (0.566).
\item \emph{The residual term is a large part of the width.} Targeting $\Bt$ instead
of $B$ narrows CLR by 0.06--0.10 in the $3\times3$ design and 0.06--0.19 in the
$2\times2$ M1 designs. The Bhatia--Davis allowance narrows CLR for $B$ from 0.086 to
0.057 in the boundary design C; where the plausible rates include 0.5 it equals
Popoviciu's and is up to 0.015 wider, because part of $\alpha$ is spent on $\bar m$.
The projection interval for $B$ with the proved allowance of
Proposition~\ref{prop:bd} differs in width by at most 0.002 from the version with a
CLR input, and its coverage is the same in every design.
\item \emph{Rank.} The numerical rank of $\hat\Sigma$ equals $\dim\beta$ in every
replication of every $2\times2$ design ($\dim\beta$ from 3 to 401). The $3\times3$
fit on all three columns has rank 12 of 18 because of the exact adding-up identity;
on the $C-1$ free columns it has full rank 12.
\end{enumerate}

\paragraph{District shocks.} Table~\ref{tab:shocks} compares M0 and M1 in a
$2\times2$ design with 20 districts of 50 units. The district shifts are added to
both groups' rates and depend on the district's mean composition $m_d$: none (S0);
affine in $m_d$ plus noise (S1); or a common shock $0.25\cos(2\pi m_d)$, which is not
affine (S2). With a non-affine shock every M0 interval fails, mostly by being empty,
while M1 covers in every replication and the pooling test always rejects. With
affine shifts the pooling test rejects every time, because M0 is false for the
observed design (Proposition~\ref{prop:tests}(c)), yet M0's intervals still cover:
shifts affine in composition change the pooled conditional mean almost only
through terms the pooled model absorbs (Proposition~\ref{prop:D4}(b)). A rejection
of the pooling test says that M0 is wrong, not that its intervals are.

\begin{table}[t]
\centering\footnotesize
\setlength{\tabcolsep}{3pt}
\begin{tabular}{@{}llccccc@{}}
\toprule
Shifts & Fit & Plug-in & CLR, $\Bt$ & Projection, $\Bt$ & Projection, $B$ & Pooling test \\
\midrule
S0 none & M0 & 1 / 0.000 & 1 / 0.329 & 1 / 0.336 & 1 / 0.436 & \\
 & M1 & 0.999 / 0.000 & 1 / 0.468 & 1 / 0.704 & 1 / 0.808 & size 0.090 ($\chi^2$) \\
S1 affine & M0 & 1 / 0.000 & 1 / 0.335 & 1 / 0.344 & 1 / 0.444 & \\
 & M1 & 0.938 / 0.000 & 1 / 0.437 & 1 / 0.672 & 1 / 0.773 & rejects 1.000 \\
S2 non-affine & M0 & 0.002 / 0.987 & 0.009 (0.963) & 0.013 (0.950) & 0.206 (0.478) & \\
 & M1 & 0.997 / 0.002 & 1 / 0.328 & 1 / 0.509 & 1 / 0.611 & rejects 1.000 \\
\bottomrule
\end{tabular}
\caption{District shocks, M0 against M1: $2\times2$, 20 districts $\times$ 50 units, 1000 replications. Coverage / median width; for S2 under M0, coverage with the share of empty intervals in parentheses. The plug-in column gives coverage / share empty. The pooling-test column gives the $\chi^2$ rejection rate at 5\%.}\label{tab:shocks}
\end{table}

\paragraph{Calibration of the tests.} In design S0 the pooling test has 38
restrictions. Its HC3 Wald statistic has mean 39.7 and variance 100, against 38 and
76 for $\chi^2_{38}$, and the $\chi^2$ test rejects 0.090 (standard error 0.009) at
nominal 5\% in the 1000-replication study above. A separate 2000-replication run
with the restricted wild bootstrap ($B=399$) gives:

\begin{center}\small
\begin{tabular}{@{}lccccc@{}}
\toprule
Test & Restrictions & Replications & $\chi^2$ & F & Wild bootstrap \\
\midrule
Pooling (M0 against M1) & 38 & 2000 & 0.078 (0.006) & 0.072 (0.006) & 0.043 (0.005) \\
Curvature (M1 against M2) & 19 & 1000 & 0.057 (0.007) & 0.055 (0.007) & 0.054 (0.007) \\
\bottomrule
\end{tabular}
\end{center}

The two $\chi^2$ figures for the pooling test (0.090 and 0.078) come from different
runs and differ by about one combined standard error. The bootstrap is close to
nominal in this design; that is a simulation finding.

\subsection{A benchmark with known truth: New Zealand elections}

New Zealand publishes, for each electorate, the exact cross-tabulation of the party
vote by the electorate (candidate) vote, together with polling-place totals of both.
This gives a benchmark in which the truth is known. Units are polling places; rows
are the party vote (Labour, National, other); columns are the electorate vote
(Labour candidate, National candidate, other); seven general elections, 2002--2020.
Intervals are for the actual rate, use the Bhatia--Davis allowance, and are
intersected with Duncan--Davis. The benchmark was run with version 0.2.0 of the
package, before the later changes to the tests and to the variance bound of the
projection interval for $B$ (Proposition~\ref{prop:bd}); those changes can move the
projection and test columns slightly.

\paragraph{One table per electorate.} There are 492 electorate tables with a median
of 64 polling places each, and 4302 cells. The pooled model's CLR and projection
intervals covered every cell, at about Duncan--Davis width (median width ratio
1.00; CLR median width 0.833 against 0.843 for DD). Adding the ordering belief that
each party's voters are at least as likely as the other party's voters to vote for
that party's candidate also covered every cell. Goodman's constancy (L6) gave
intervals with median width 0.415 that covered 0.939 of cells and missed the truth
with a non-empty interval in 2.2\%. The plug-in estimate of the M0 identified set was
empty in 48\% of cells. With about 64 units per table, the valid intervals are about
as wide as Duncan--Davis.

\paragraph{Electorates as districts.} Pooling each election's 62--69 electorates
(3398--6550 polling places) and fitting M1 with electorates as districts, the data
reject the model: the M1 plug-in estimate is empty in all seven elections, and the
curvature and pooling tests have $p\approx0$ (the pooling test has 366--408
restrictions). The emptiness test itself rejects the pooled model in two elections
and M1 in none. Over the 63 cells (seven elections, nine cells), the M1 CLR
intervals had crossing ends (or ends outside the logical range) in 21\% of cells and missed the truth with a non-empty
interval in 8\% (5 of 63); the projection interval never missed, but it was as wide
as Duncan--Davis. Where the model is wrong, the tests flag it, and a non-empty CLR
interval is not reassurance.

\subsection{The sign module}

\paragraph{How often signs add information.} In 2000 random $2\times2$ LCE designs
with true signs imposed, the within sign strictly tightened JKST's interval in 83\%
of designs (median width ratio 0.70), the between sign in 89\% (0.67), and CR, with
unknown sign, in 80\% of the 1060 designs where CR holds (0.69). In a $3\times2$
design with true-sign rows on both associations, a row was implied in 8\% of
target-row pairs, cut the set but not the target in 28\%, and tightened the target in
64\% (Proposition~\ref{prop:R2}).

\paragraph{CR at the edge.} Table~\ref{tab:cr} uses a $2\times2$ design with 150
units, 600 replications, Beta rates with linear means (so LCE and CR hold) and the
conservative interval for the gap $\tilde D$. In the first row the truth lies on the
edge of the CR band, so $\delta^*_W=0$.

\begin{table}[t]
\centering\footnotesize
\begin{tabular}{@{}lcccc@{}}
\toprule
Design & JKST alone & CR slab & CR branch union & CR plug-in bracket \\
\midrule
Truth on the CR edge, $\ER^*=0.15$ & 1.000 / 1.106 & 0.992 / 0.359 & 0.992 / 0.357 & \textbf{0.917} / 0.256 \\
Truth on the edge, $\ER^*=0.005$ & 1.000 / 1.223 & 1.000 / 0.242 & 1.000 / 0.242 & 1.000 / 0.147 \\
Interior truth & 1.000 / 1.232 & 1.000 / 0.480 & 1.000 / 0.480 & 1.000 / 0.388 \\
\bottomrule
\end{tabular}
\caption{Contextual reinforcement in $2\times2$: coverage / median length of conservative intervals for the gap (600 replications). The branch union is measured by its occupied length. No interval was empty.}\label{tab:cr}
\end{table}

The slab cuts the median length from 1.106 to 0.359 at the edge and keeps coverage.
The plug-in bracket, which imposes $\hull\{\gamma\widehat\ER,\widehat\ER\}$ and then
adds the target's own inflation, covers 0.917, consistent with the limit of
Theorem~\ref{thm:R10}(e) (at most 0.9335 in the three-point design). The literal
intersection of the fitted bracket with a finished JKST interval covered 304 of 600
(51\%), in line with its limit of at most one half. The package's default CLR branch
union, in a separate run of 200 replications of the edge design, covered 0.990 with
median width 0.230; that interval is algorithmic.

\paragraph{A false belief.} With the same generator and means $m_0=.40-sx$ and
$m_1=.50+s-sx$, the aggregate regression is $.40+.10x$ for every $s$, the between
sign is positive throughout, and for $s>0$ both contextual slopes are negative, so
CR is false.

\begin{center}\small
\begin{tabular}{@{}rrccccc@{}}
\toprule
Units & $s$ & True $\tilde D$ & JKST covers & CR slab covers & JKST width & CR width \\
\midrule
150 & 0 & .100 & 400/400 & 399/400 & 1.262 & .306 \\
150 & .08 & .170 & 400/400 & 336/400 & 1.256 & .304 \\
150 & .25 & .318 & 400/400 & 2/400 & 1.254 & .303 \\
600 & 0 & .100 & 400/400 & 398/400 & 1.074 & .194 \\
600 & .08 & .169 & 400/400 & 110/400 & 1.069 & .196 \\
600 & .25 & .317 & 400/400 & 0/400 & 1.066 & .195 \\
\bottomrule
\end{tabular}
\end{center}

No interval was empty. A false restriction gives narrow, non-empty and reproducibly
wrong intervals: precision makes the wrong answer more reproducible, not more
credible. This is the price stated in Theorem~\ref{thm:R10}(d).

\subsection{A worked example}\label{sec:example}

We illustrate the package on the \code{voter\_vax} data distributed with the
MonotoneEI package of EGGHM (and included in \pkg{eiPartialID}'s tests): 2059 US
counties, with $X$ the Republican share of the presidential vote, $T$ the share
vaccinated against COVID-19, and $N$ the total votes cast. Rows are Republican and
other voters; the column is ``vaccinated''; the target is the gap $D$ between the two
groups' vaccination rates. The example illustrates the mechanics. It is not a
substantive finding: the vaccination share refers to the county population and the
group shares to voters, so the closed-population assumption A1 holds only
approximately.

Fitting the pooled rung with size weights,
\begin{center}
\code{eiBounds(x = prop\_republican, t = prop\_covid\_vax, n = total\_votes, weights = "n",}\\
\code{targets = list(D = matrix(c(1, -1, 0, 0), 2, 2)), ci = "clr")}
\end{center}
gives the results in Table~\ref{tab:vv} (seed 1). The estimated identified set for
the gap under LCE alone is $[-0.486,-0.001]$: Republican voters' vaccination rate is
lower, by between 0 and 49 points. The projection interval, which carries the
asymptotic guarantee for the pooled model, is $[-0.607,0.080]$, and the CLR interval
is slightly narrower. The actual-rate intervals are much wider, because the residual
allowance covers each county's local luck.

\begin{table}[t]
\centering\small
\begin{tabular}{@{}lcccc@{}}
\toprule
& \multicolumn{2}{c}{M0 alone} & \multicolumn{2}{c}{M0 + CR} \\
\cmidrule(lr){2-3}\cmidrule(lr){4-5}
& Gap $D$ & Republican rate & Gap $D$ & Republican rate \\
\midrule
Estimate (identified set) & $[-0.486,-0.001]$ & $[0.345,0.597]$ & $[-0.470,-0.050]$ & $[0.354,0.571]$ \\
CLR, $\Bt$ (algorithmic) & $[-0.571,\ \ 0.056]$ & $[0.301,0.623]$ & $[-0.525,-0.044]$ & $[0.330,0.580]$ \\
Projection, $\Bt$ & $[-0.607,\ \ 0.080]$ & $[0.282,0.634]$ & $[-0.542,-0.043]$ & $[0.323,0.582]$ \\
Conservative, $\Bt$ & $[-0.599,\ \ 0.112]$ & $[0.291,0.651]$ & $[-0.530,\ \ 0.011]$ & $[0.326,0.600]$ \\
CLR, $B$ (algorithmic) & $[-0.750,\ \ 0.233]$ & $[0.229,0.694]$ & $[-0.702,\ \ 0.131]$ & $[0.259,0.650]$ \\
Projection, $B$ & $[-0.763,\ \ 0.257]$ & $[0.210,0.705]$ & $[-0.719,\ \ 0.133]$ & $[0.252,0.652]$ \\
Duncan--Davis & $[-0.763,\ \ 0.547]$ & $[0.202,0.881]$ & $[-0.763,\ \ 0.547]$ & $[0.202,0.881]$ \\
\bottomrule
\end{tabular}
\caption{The \code{voter\_vax} example: 2059 counties, pooled rung, size-weighted regression, 95\% intervals. The conservative row comes from a separate fit with \code{ci = "conservative"}.}\label{tab:vv}
\end{table}

\paragraph{Descriptive quantities.} The Republican share has mean $\mu=0.482$ and
between-county variance $0.0267$, so $\gamma=0.107$. The regression slope is
$\ER=-0.470$ and the neighbourhood gap is $\NM=\gamma\ER=-0.050$. By
Proposition~\ref{prop:G2}, each sign belief cuts the identified set at one of these
two numbers:
\begin{center}\small
\begin{tabular}{@{}ll@{}}
\toprule
Belief & Estimated set for $D$ \\
\midrule
$\delta_W\ge0$ (vaccination rises with the Republican share, on average) & $[-0.486,-0.470]$ \\
$\delta_W\le0$ & $[-0.470,-0.001]$ \\
$\delta_B\ge0$ (Republicans more likely to be vaccinated within a county) & $[-0.050,-0.001]$ \\
$\delta_B\le0$ & $[-0.486,-0.050]$ \\
CR (the two associations share a sign) & $[-0.470,-0.050]$ \\
\bottomrule
\end{tabular}
\end{center}
Under the rung alone neither association's sign is settled: the estimated set for
$\delta_B$ is $[-0.488,0.055]$ and for $\delta_W$ is $[-0.525,0.018]$.
\code{eiRedundancy()} reports that CR tightens the estimated set for $D$ by 0.066,
whether added first or last. Since $\ER<0$, only the negative branch of CR is
non-empty (Proposition~\ref{prop:tworowsign}(d)), and the CR set is the band between
$\ER$ and $\NM$. With CR, the projection interval for $\Bt$, $[-0.542,-0.043]$,
excludes zero: under LCE and CR, Republican voters' model-implied vaccination rate
is lower, with an asymptotic guarantee for the pooled model
(Theorem~\ref{thm:R10}). The conclusion is only as credible as CR.

\paragraph{A joint rejection.} Imposing both $\delta_W\ge0$ and $\delta_B\ge0$ asks
for $D\le\ER=-0.470$ and $D\ge\NM=-0.050$ at once, which is impossible. The estimated
set is empty, and the emptiness test rejects the combination of the rung and the two
beliefs, while the rung alone is not rejected. The package reports exactly this: the
empty estimate is a sample diagnostic, the emptiness test is the test, the rejection
is of the combination and does not say which part fails, and each belief on its own
gives a non-empty set (Theorem~\ref{thm:R3}(e)). The signed associations have
effective numbers of counties of 271 ($\delta_B$) and 188 ($\delta_W$), so no single
county dominates them.

% Section 9
\section{Open questions}\label{sec:open}

The proved results are stated in Sections~\ref{sec:rxc}--\ref{sec:inference}.
This section collects what is not proved: the conjectures and algorithmic
procedures that appear in the paper, and the open problems.

\subsection{Conjectures and algorithmic procedures}

\begin{itemize}
\item \textbf{CLR as the default interval} (Conjecture~\ref{conj:clr}). The
package's default confidence interval collects candidate endpoints by simulation
and computes its selection threshold over the collected set. It is an algorithm.
The reduction to the cited full-candidate CLR theorem is proved
(Lemma~\ref{lem:reduction}); the recovery of rarely optimal vertices, the
anti-conservative effect of computing the threshold over a subset, a growing
number of candidates, and uniform validity are open. Its coverage was at least
0.994 wherever Slater's condition held, and nominal at a point-identified design,
but simulation cannot make it a theorem. The variance-adaptive allowance computed
from a CLR interval, and the CLR branch union for contextual reinforcement,
inherit this status. Proving the conjecture would make the default a theorem.
\item \textbf{The sharp N2 set} (Conjecture~\ref{conj:N2}). Sufficiency of the
selection-expectation characterisation on a finite support is shown; necessity in
general is not.
\item \textbf{The balance point} (Conjecture~\ref{conj:C1}). That $x^*$ lies in the
composition hull in $2\times2$, which would complete the chain from signed slopes to
the within sign.
\item \textbf{Shared and district-specific effects in regime R3}: the asymptotic
results with a high-dimensional Gaussian approximation, under a rate condition on
the number of districts relative to their size.
\item \textbf{The wild-bootstrap calibration} of the pooling and curvature tests
is supported by simulation; its validity for a growing number of restrictions is
not proved.
\item \textbf{Spanning trees of pairwise beliefs}: that every spanning tree, not
only a star, reaches all directions of $K$ (checked for chains), and a related
conjecture that single-index contextual effects have dimension $C-1$.
\end{itemize}

\subsection{Open problems}

\begin{itemize}
\item \textbf{Sign rows under district rungs.} Sharpness and inference for the sign
and reinforcement rows under M1, M1($\delta$), M2 and L4. The rows are valid linear
rows, and the bound with them is an outer set, but neither its sharpness for the
regression nor the coverage of its intervals is proved. Until then the grouped
sign module is experimental and its intervals have no coverage theorem.
\item \textbf{Integer unit tables.} The integer--continuous gap beyond the
$2\times2$ case with margins only (Proposition~\ref{prop:integer}): with unit-level
restrictions, several columns or three or more groups.
\item \textbf{Inference on width gains.} Whether a belief tightens a target is
decided exactly at the population (Proposition~\ref{prop:R2}) and estimated at the
point estimate in samples. Inference on the width change needs joint treatment of
the base and augmented endpoints and of any estimated threshold; an interval for
either endpoint alone is not enough.
\item \textbf{Informative finite-sample inference with many small groups.}
Theorem~\ref{thm:finite} is valid in regime R2 but uninformative. Candidates are a
direct bound on the endpoint rather than on every constraint, variance-adaptive
(empirical-Bernstein) radii, a Bhatia--Davis version of the residual radius, and
uneven splits of the error level. A hierarchical model for district shifts would
change the model.
\item \textbf{Uniform limit theory under changing active sets.}
Proposition~\ref{prop:triangular} covers convergent localised slacks; a theory
uniform in the design sequence, when candidates stay within $O(n^{-1/2})$ of binding
without converging, is open. So are uniform validity near point identification for
the CLR interval (a Stoye-type adaptation) and the rate and limit of the endpoint
estimator when Slater's condition fails.
\item \textbf{The hull error.} Its rate, and inference for the population-hull
endpoint.
\item \textbf{Other targets and assumptions.} Inference for the superpopulation
target $B^*$; dependence within districts beyond the district effect; a joint
relaxation of the closed-population assumption for turnover; coverage of the
plug-in under misspecification; the characterisation of strict gains from joint
fitting and from shared district effects when $R\ge3$; refutability of sign,
ordering and local beliefs when $R\ge3$; an outer N2 LP with an estimated
regression; and a consistent series bound that would tighten $\ThO$ towards $\ThX$,
which conditions on $X$ only and so would not be sharp when size is informative.
\end{itemize}

\subsection{Future work}

Taking election periods as districts turns the shared-effects rung into the
assumption that contextual effects are unchanged between elections, which opens a
route to bounds on changes in ecological rates and to difference-in-differences
designs; that extension is left for future work.


\bibliographystyle{plainnat}
\bibliography{references}

\end{document}
